\documentclass[11pt,a4paper,oneside]{book}
\usepackage[utf8]{inputenc}
\usepackage[T1]{fontenc}
\usepackage{lmodern}
\usepackage[english]{babel}
\usepackage{amsmath,amssymb,amsthm,mathtools}
\usepackage{booktabs}
\usepackage{graphicx}
\usepackage{xcolor}
\usepackage{tikz}
\usetikzlibrary{arrows.meta,positioning,calc,shapes,shapes.geometric,decorations.pathreplacing,fit,backgrounds,trees}
\usepackage{listings}
\usepackage{hyperref}
\hypersetup{colorlinks=true,linkcolor=blue!60!black,citecolor=green!50!black,urlcolor=blue!60!black}
\usepackage{enumitem}
\usepackage{geometry}
\geometry{a4paper, top=1.3cm, bottom=1.3cm, left=1.2cm, right=1.2cm, includeheadfoot}
\usepackage{setspace}
\setstretch{1.20}
\setlength{\parskip}{0.70em}
\setlength{\parindent}{1.0em}
\setlength{\abovedisplayskip}{10pt plus 2pt minus 2pt}
\setlength{\belowdisplayskip}{10pt plus 2pt minus 2pt}
\setlength{\abovedisplayshortskip}{7pt plus 2pt}
\setlength{\belowdisplayshortskip}{7pt plus 2pt}
\setlist{topsep=0.55em, partopsep=0.25em, itemsep=0.45em, parsep=0.25em}
\setlength{\textfloatsep}{10pt plus 2pt minus 2pt}
\setlength{\intextsep}{8pt plus 2pt minus 2pt}

\theoremstyle{definition}
\newtheorem{definition}{Definition}[chapter]
\newtheorem{theorem}{Theorem}[chapter]
\newtheorem{lemma}{Lemma}[chapter]
\newtheorem{corollary}{Corollary}[chapter]
\newtheorem{example}{Example}[chapter]
\newtheorem{exercise}{Exercise}[chapter]
\theoremstyle{remark}
\newtheorem{remark}{Remark}[chapter]

\lstset{basicstyle=\ttfamily\footnotesize,breaklines=true,frame=single,numbers=left,numberstyle=\tiny\color{gray},keywordstyle=\color{blue!70!black},commentstyle=\color{green!50!black},stringstyle=\color{orange!70!black},showstringspaces=false,tabsize=2, xleftmargin=1.0em, xrightmargin=0.5em, aboveskip=0.6em, belowskip=0.6em}

\title{SC1004 Linear Algebra for Computing\\NTU BSc Computer Science --- Comprehensive Textbook (0--100)}
\author{NTU CS Curriculum Textbook Series\\College of Computing and Data Science}
\date{2026-08-23 --- Generated for bumbsclaw/ntu-cs-textbooks}

\begin{document}
\frontmatter
\maketitle
\tableofcontents
\mainmatter
% SC1004 Chapter 1: Vectors & Geometry (0->100 ground-up)
\chapter{Vectors and Geometry}
\label{ch:vectors}

\begin{quote}
\emph{``A vector is not just a list of numbers --- it is an arrow that knows where to push.''} This chapter gives that arrow a home: from points on a map to $n$-dimensional space, with pictures before formulas.
\end{quote}

\noindent\fbox{\parbox{\dimexpr\linewidth-2\fboxsep-2\fboxrule}{%
\textbf{From zero: no prior linear algebra assumed.} We start with a dot on paper, then an arrow, then coordinates. Every notion --- length, angle, line, plane --- is drawn before it is defined. By the end you will do geometry with algebra.}}
\vspace{0.4em}

\section*{Learning Outcomes}
After this chapter you will be able to:
\begin{enumerate}[label=(L\arabic*)]
    \item represent vectors geometrically and algebraically in $\mathbb{R}^2$, $\mathbb{R}^3$, and $\mathbb{R}^n$;
    \item compute dot products, norms, and angles and interpret them geometrically;
    \item write parametric equations of lines and planes;
    \item use vector geometry to solve distance and projection problems;
    \item translate between pictures, coordinates, and code (NumPy).
\end{enumerate}

% ============================================================
\section{What Is a Vector?}
\label{sec:what-is-vector}
Intuition: stand on a map. A \emph{point} says where you are. A \emph{vector} says how to move: ``3 east, 2 north.'' That move has a length and a direction, but no fixed location --- sliding the arrow keeps the same vector.

\begin{figure}[h]
\centering
\begin{tikzpicture}[scale=0.85,every node/.style={font=\scriptsize}]
  \draw[->,thick] (-0.5,0)--(4.2,0) node[right] {$x$};
  \draw[->,thick] (0,-0.5)--(0,3.2) node[above] {$y$};
  \fill[blue!60!black] (0.6,0.6) circle (2pt) node[below left] {$P$};
  \fill[red!70!black] (3.2,2.3) circle (2pt) node[above right] {$Q$};
  \draw[->,very thick,blue!60!black] (0.6,0.6)--(3.2,2.3) node[midway,above,sloped] {$\vec{v}=\overrightarrow{PQ}$};
  \draw[->,very thick,orange!75!black,dashed] (0.4,0.2)--(3.0,1.9) node[midway,below,sloped] {same $\vec{v}$};
  \draw[dashed,gray] (0.6,0.6)--(3.2,0.6)--(3.2,2.3);
  \node[font=\tiny,fill=white] at (1.9,0.35) {$\Delta x=2.6$};
  \node[font=\tiny,fill=white,rotate=90] at (3.45,1.45) {$\Delta y=1.7$};
\end{tikzpicture}
\caption{A vector as a free arrow. $\overrightarrow{PQ}$ and its translate have the same components $(\Delta x,\Delta y)$.}
\label{fig:free-vector}
\end{figure}

\begin{definition}[Vector in $\mathbb{R}^n$]
A \emph{vector} $\mathbf{v}\in\mathbb{R}^n$ is an ordered $n$-tuple
\[
  \mathbf{v}=\begin{bmatrix}v_1\\v_2\\\vdots\\v_n\end{bmatrix},\quad v_i\in\mathbb{R}.
\]
We write $\mathbf{0}$ for the zero vector and $-\mathbf{v}$ for the opposite arrow. Two arrows represent the same vector iff their component differences agree.
\end{definition}

Addition is tip-to-tail; scalar multiplication stretches.

\begin{figure}[h]
\centering
\begin{tikzpicture}[scale=0.82,every node/.style={font=\scriptsize}]
  \begin{scope}
    \draw[->,thick] (-0.3,0)--(3.0,0); \draw[->,thick] (0,-0.3)--(0,2.5);
    \draw[->,very thick,blue!65!black] (0,0)--(1.8,0.6) node[midway,below] {$\mathbf{u}$};
    \draw[->,very thick,red!65!black] (1.8,0.6)--(2.8,2.0) node[midway,right] {$\mathbf{v}$};
    \draw[->,very thick,violet!70!black] (0,0)--(2.8,2.0) node[midway,above left] {$\mathbf{u}+\mathbf{v}$};
    \node[font=\footnotesize\bfseries] at (1.4,-0.7) {Tip-to-tail sum};
  \end{scope}
  \begin{scope}[xshift=4.5cm]
    \draw[->,thick] (-0.3,0)--(3.0,0); \draw[->,thick] (0,-0.3)--(0,2.5);
    \draw[->,very thick,blue!65!black] (0,0)--(1.5,1.0) node[midway,above left] {$\mathbf{u}$};
    \draw[->,very thick,orange!75!black] (0,0)--(3.0,2.0) node[midway,below right] {$2\mathbf{u}$};
    \draw[->,very thick,green!55!black] (0,0)--(-0.75,-0.5) node[midway,below] {$-0.5\mathbf{u}$};
    \node[font=\footnotesize\bfseries] at (1.2,-0.7) {Scalar multiples};
  \end{scope}
\end{tikzpicture}
\caption{Left: $\mathbf{u}+\mathbf{v}$ (parallelogram rule). Right: $c\mathbf{u}$ scales length by $|c|$, flips if $c<0$.}
\label{fig:add-scale}
\end{figure}

Algebraically,
\[
  \mathbf{u}+\mathbf{v}=\begin{bmatrix}u_1+v_1\\\vdots\\u_n+v_n\end{bmatrix},\quad
  c\mathbf{u}=\begin{bmatrix}cu_1\\\vdots\\cu_n\end{bmatrix}.
\]
Properties: commutativity, associativity, distributivity --- all inherited from $\mathbb{R}$.

\section{Length, Dot Product, and Angle}
\label{sec:dot}
How long is an arrow? Pythagoras in $n$ dimensions.

\begin{definition}[Norm and dot product]
For $\mathbf{u},\mathbf{v}\in\mathbb{R}^n$,
\[
  \mathbf{u}\cdot\mathbf{v}=u_1v_1+\cdots+u_nv_n,\qquad
  \|\mathbf{u}\|=\sqrt{\mathbf{u}\cdot\mathbf{u}}=\sqrt{u_1^2+\cdots+u_n^2}.
\]
Distance between tips: $\|\mathbf{u}-\mathbf{v}\|$.
\end{definition}

\begin{theorem}[Cauchy--Schwarz and angle]
For any $\mathbf{u},\mathbf{v}$,
\[
  |\mathbf{u}\cdot\mathbf{v}|\le\|\mathbf{u}\|\,\|\mathbf{v}\|,
\]
with equality iff one is a scalar multiple of the other. If both nonzero, there is a unique $\theta\in[0,\pi]$ with
\[
  \mathbf{u}\cdot\mathbf{v}=\|\mathbf{u}\|\,\|\mathbf{v}\|\cos\theta,
\]
and $\theta$ is the geometric angle between them.
\end{theorem}
\begin{proof}[Sketch]
Consider $\|\mathbf{u}-t\mathbf{v}\|^2\ge0$ for all $t\in\mathbb{R}$; its discriminant must be $\le0$, giving Cauchy--Schwarz. Define $\theta$ via $\cos\theta=(\mathbf{u}\cdot\mathbf{v})/(\|\mathbf{u}\|\|\mathbf{v}\|)$.
\end{proof}

Corollary: $\mathbf{u}\perp\mathbf{v}$ (orthogonal) iff $\mathbf{u}\cdot\mathbf{v}=0$.

\begin{figure}[h]
\centering
\begin{tikzpicture}[scale=0.85,every node/.style={font=\scriptsize}]
  \draw[->,thick] (-0.3,0)--(4.0,0) node[right] {$x$};
  \draw[->,thick] (0,-0.3)--(0,3.0) node[above] {$y$};
  \draw[->,very thick,blue!65!black] (0,0)--(3.0,0.7) node[midway,below] {$\mathbf{u}$};
  \draw[->,very thick,red!65!black] (0,0)--(1.2,2.4) node[midway,left] {$\mathbf{v}$};
  \draw[dashed,gray] (1.2,2.4)--(2.15,0.78);
  \fill[orange!30] (0,0)--(1.0,0.23)--(0.92,0.45)--cycle;
  \node at (0.55,0.18) {$\theta$};
  \draw[->,thick,violet!70!black] (0,0)--(2.15,0.78) node[midway,below] {$\operatorname{proj}_{\mathbf{u}}\mathbf{v}$};
  \node[font=\tiny,fill=white] at (1.7,1.65) {$\mathbf{v}-\operatorname{proj}$};
\end{tikzpicture}
\caption{Angle $\theta$ via the dot product; projection of $\mathbf{v}$ onto $\mathbf{u}$ is the shadow along $\mathbf{u}$.}
\label{fig:angle-proj}
\end{figure}

\emph{Projection:} the shadow of $\mathbf{v}$ on $\mathbf{u}\neq\mathbf{0}$ is
\[
  \operatorname{proj}_{\mathbf{u}}\mathbf{v}=\frac{\mathbf{v}\cdot\mathbf{u}}{\mathbf{u}\cdot\mathbf{u}}\,\mathbf{u},\qquad
  \|\operatorname{proj}_{\mathbf{u}}\mathbf{v}\|= \|\mathbf{v}\|\,|\cos\theta|.
\]
Useful everywhere from graphics (shadows) to ML (closest point).

\begin{example}[Cosine similarity]
In text retrieval, documents are vectors of word counts. $\cos\theta=(\mathbf{a}\cdot\mathbf{b})/(\|\mathbf{a}\|\|\mathbf{b}\|)$ near $1$ means similar topics --- the same formula as geometry.
\end{example}

\section{Lines and Planes}
\label{sec:lines-planes}
A line is a point plus a direction. A plane is a point plus two independent directions.

\begin{align}
  \text{Line: } &\mathbf{x}=\mathbf{p}+t\mathbf{d},\quad t\in\mathbb{R},\\
  \text{Plane in $\mathbb{R}^3$: } &\mathbf{x}=\mathbf{p}+s\mathbf{d}_1+t\mathbf{d}_2,\quad s,t\in\mathbb{R},\\
  \text{Plane (normal form): } &\mathbf{n}\cdot(\mathbf{x}-\mathbf{p})=0\iff \mathbf{n}\cdot\mathbf{x}=d.
\end{align}
Here $\mathbf{d}$ is a direction vector and $\mathbf{n}$ is a \emph{normal} (perpendicular) vector. In $\mathbb{R}^3$, if $\mathbf{d}_1,\mathbf{d}_2$ span the plane then $\mathbf{n}=\mathbf{d}_1\times\mathbf{d}_2$ (cross product).

\begin{example}[Distance from point to line]
Let $L:\mathbf{p}+t\mathbf{d}$. For $\mathbf{q}$, let $\mathbf{w}=\mathbf{q}-\mathbf{p}$. Then
\[
  \text{dist}(\mathbf{q},L)=\left\|\mathbf{w}-\operatorname{proj}_{\mathbf{d}}\mathbf{w}\right\|
  =\frac{\|\mathbf{w}\times\mathbf{d}\|}{\|\mathbf{d}\|}\quad(\text{in }\mathbb{R}^3).
\]
Subtract the along-line shadow; what remains is the perpendicular distance.
\end{example}

\begin{example}[Python --- vectors as arrays]
\begin{lstlisting}[language=Python]
import numpy as np
u = np.array([3.0, 1.0])
v = np.array([1.0, 2.0])
dot = u @ v                      # 5.0
norm_u = np.linalg.norm(u)       # sqrt(10)
cos_theta = dot/(norm_u*np.linalg.norm(v))
proj = (dot/(u@u))*u             # projection of v onto u? careful: swap
proj_v_on_u = (v@u)/(u@u) * u
print(cos_theta, proj_v_on_u)
\end{lstlisting}
\end{example}

% ============================================================

% --- Additional depth: Cross product and applications ---
\section{Cross Product and Volume}
\label{sec:cross}
In $\mathbb{R}^3$, the \emph{cross product} $\mathbf{u}\times\mathbf{v}$ is the vector orthogonal to both $\mathbf{u},\mathbf{v}$ with magnitude $\|\mathbf{u}\|\|\mathbf{v}\|\sin\theta$ (area of the parallelogram) and direction by the right-hand rule.
\[
  \mathbf{u}\times\mathbf{v}=\begin{bmatrix}u_2v_3-u_3v_2\\u_3v_1-u_1v_3\\u_1v_2-u_2v_1\end{bmatrix},\quad
  \|\mathbf{u}\times\mathbf{v}\|=\|\mathbf{u}\|\|\mathbf{v}\|\sin\theta.
\]
Properties: $\mathbf{u}\times\mathbf{v}=-(\mathbf{v}\times\mathbf{u})$, $\mathbf{u}\times\mathbf{u}=\mathbf{0}$, bilinear. The scalar triple product $\mathbf{u}\cdot(\mathbf{v}\times\mathbf{w})=\det[\mathbf{u}\ \mathbf{v}\ \mathbf{w}]$ is the signed volume.

\begin{example}[Normal to a plane]
Points $P(1,0,0),Q(0,1,0),R(0,0,1)$: $\overrightarrow{PQ}=[-1,1,0]^T$, $\overrightarrow{PR}=[-1,0,1]^T$, so $\mathbf{n}=\overrightarrow{PQ}\times\overrightarrow{PR}=[1,1,1]^T$. Plane: $x+y+z=1$, matching $\mathbf{n}\cdot(\mathbf{x}-P)=0$.
\end{example}

\section{Inequalities and Geometry}
\label{sec:inequalities}
Two inequalities govern all distance arguments:
\begin{theorem}[Triangle and Cauchy--Schwarz]
For $\mathbf{u},\mathbf{v}\in\mathbb{R}^n$: (i) $\|\mathbf{u}+\mathbf{v}\|\le\|\mathbf{u}\|+\|\mathbf{v}\|$ (triangle), equality when $\mathbf{u},\mathbf{v}$ same direction; (ii) $|\mathbf{u}\cdot\mathbf{v}|\le\|\mathbf{u}\|\|\mathbf{v}\|$ with equality iff dependent.
\end{theorem}
Triangle follows from $\|\mathbf{u}+\mathbf{v}\|^2=\|\mathbf{u}\|^2+2\mathbf{u}\cdot\mathbf{v}+\|\mathbf{v}\|^2\le(\|\mathbf{u}\|+\|\mathbf{v}\|)^2$ via Cauchy--Schwarz. These justify calling $\|\cdot\|$ a norm and $d(\mathbf{u},\mathbf{v})=\|\mathbf{u}-\mathbf{v}\|$ a metric.


\section{Chapter Notes}
Vectors go back to Grassmann and Gibbs; modern computing treats them as arrays. The dot product as $\|\mathbf{u}\|\|\mathbf{v}\|\cos\theta$ unifies algebra and geometry --- the bridge used in graphics, robotics, and embeddings.

% ============================================================
\section{Exercises}
\label{sec:ch01-exercises}
\begin{enumerate}[label=E1.\arabic*, leftmargin=*, itemsep=0.8em]
\item \textbf{Arithmetic.} Let $\mathbf{u}=[2,-1,3]^T$, $\mathbf{v}=[0,2,-1]^T$. Compute $\mathbf{u}+\mathbf{v}$, $3\mathbf{u}-2\mathbf{v}$, $\mathbf{u}\cdot\mathbf{v}$, $\|\mathbf{u}\|$, and the angle $\theta$ between $\mathbf{u}$ and $\mathbf{v}$ (to $1^\circ$).
\item \textbf{Orthogonality.} Find all $c$ such that $[c,2,1]^T$ is orthogonal to $[1,c,-2]^T$. Give a geometric interpretation of your answer.
\item \textbf{Line and distance.} Let $L:\mathbf{p}=[1,0,2]^T$, $\mathbf{d}=[1,1,1]^T$. Write the parametric equation, find the point on $L$ closest to $\mathbf{q}=[2,3,0]^T$, and compute $\text{dist}(\mathbf{q},L)$.
\item \textbf{Projection proof.} Prove $\mathbf{v}-\operatorname{proj}_{\mathbf{u}}\mathbf{v}$ is orthogonal to $\mathbf{u}$. Use this to derive the formula for distance from a point to a plane $\mathbf{n}\cdot\mathbf{x}=d$ in $\mathbb{R}^3$.
\item \textbf{Coding.} Write a Python function \texttt{gram\_schmidt\_2(u,v)} returning an orthonormal pair spanning the same plane as $\mathbf{u},\mathbf{v}$ (assume independent). Test on $\mathbf{u}=[1,1,0]^T$, $\mathbf{v}=[1,0,1]^T$ and verify orthonormality numerically.
\end{enumerate}
% End of Chapter 1

% SC1004 Chapter 2: Matrices (0->100 ground-up)
\chapter{Matrices}
\label{ch:matrices}

\begin{quote}
\emph{``A matrix is a spreadsheet that learned to multiply.''} We move from single arrows (vectors) to tables of numbers and discover the arithmetic that powers graphics, ML, and networks.
\end{quote}

\noindent\fbox{\parbox{\dimexpr\linewidth-2\fboxsep-2\fboxrule}{%
\textbf{From zero: tables before theorems.} A matrix is first a rectangular array --- rows and columns you already know from spreadsheets. We give it addition, scaling, multiplication, and transpose, then ask what each means for vectors. No prior matrix theory assumed.}}
\vspace{0.4em}

\section*{Learning Outcomes}
After this chapter you will be able to:
\begin{enumerate}[label=(L\arabic*)]
    \item add, scale, multiply, and transpose matrices fluently;
    \item interpret $A\mathbf{x}$ as a linear combination of columns and as a linear map;
    \item recognise identity, diagonal, symmetric, and triangular matrices;
    \item compute with block matrices and outer products;
    \item implement matrix operations in NumPy and reason about cost.
\end{enumerate}

% ============================================================
\section{Anatomy of a Matrix}
\label{sec:anatomy}
Picture a spreadsheet: $m$ rows, $n$ columns, an entry in each cell.

\begin{definition}[Matrix]
An $m\times n$ \emph{matrix} over $\mathbb{R}$ is an array $A=[a_{ij}]$ with $i=1..m$, $j=1..n$. We write $A\in\mathbb{R}^{m\times n}$. A column vector is $m\times1$; a row vector $1\times n$; a scalar $1\times1$.
\end{definition}

\begin{figure}[h]
\centering
\begin{tikzpicture}[scale=0.85,every node/.style={font=\scriptsize}]
  \foreach \r in {0,...,2} \foreach \c in {0,...,3} {
    \node[draw, thick, minimum width=0.95cm, minimum height=0.58cm, fill=blue!7] at (\c*1.05,\r*-0.68) {$a_{\the\numexpr\r+1\relax\the\numexpr\c+1\relax}$};
  }
  \foreach \c in {0,...,3} \node[font=\tiny,blue!60!black] at (\c*1.05,0.45) {col \the\numexpr\c+1\relax};
  \foreach \r in {0,...,2} \node[font=\tiny,red!60!black,anchor=east] at (-0.6,\r*-0.68) {row \the\numexpr\r+1\relax};
  \draw[decorate,decoration={brace,amplitude=4pt},thick] (-0.55,0.25)--(-0.55,-1.82) node[midway,left=5pt,font=\tiny] {$m=3$};
  \draw[decorate,decoration={brace,amplitude=4pt},thick] (-0.15,-2.05)--(3.25,-2.05) node[midway,below=4pt,font=\tiny] {$n=4$};
  \node[draw,rounded corners=3pt,fill=orange!18,font=\tiny,align=center] at (5.5, -0.4) {Shape $m\times n$\\$m$ rows, $n$ cols};
  \draw[->,thick,orange!70!black] (3.8,-0.5)--(4.55,-0.5);
\end{tikzpicture}
\caption{An $m\times n$ matrix as a table. Row $i$, column $j$ holds $a_{ij}$.}
\label{fig:matrix-anatomy}
\end{figure}

Addition and scalar multiplication are entrywise (same shape required). The zero matrix $0_{m\times n}$ and negatives are as for vectors.

\section{Matrix Multiplication}
\label{sec:mult}
This is the interesting one. Not entrywise --- it composes linear actions.

\begin{definition}[Product]
If $A\in\mathbb{R}^{m\times n}$ and $B\in\mathbb{R}^{n\times p}$, then $C=AB\in\mathbb{R}^{m\times p}$ with
\[
  c_{ij}=\sum_{k=1}^{n} a_{ik}b_{kj}\quad\text{(row $i$ of $A$ dot column $j$ of $B$).}
\]
The inner dimensions must match ($n=n$); the outer give the shape.
\end{definition}

\begin{figure}[h]
\centering
\begin{tikzpicture}[scale=0.80,every node/.style={font=\scriptsize}]
  \node[draw,thick,fill=blue!12,minimum width=2.0cm,minimum height=1.2cm] (A) at (0,0) {$A:m\times n$};
  \node[draw,thick,fill=red!12,minimum width=2.0cm,minimum height=1.2cm] (B) at (3.2,0) {$B:n\times p$};
  \node[draw,thick,fill=violet!15,minimum width=2.4cm,minimum height=1.2cm] (C) at (6.8,0) {$C=AB:m\times p$};
  \draw[->,very thick,blue!60!black] (A)--(B) node[midway,above,font=\tiny] {inner $n$ matches};
  \draw[->,very thick,violet!60!black] (1.6,0.9) to[out=45,in=135] (6.8,0.9) node[midway,above,font=\tiny] {outer $m\times p$};
  \node[font=\tiny,align=center] at (1.6,-1.1) {row $\cdot$ column\\$\sum_k a_{ik}b_{kj}$};
  \begin{scope}[yshift=-2.0cm,xshift=1.2cm]
    \draw[fill=blue!15] (0,0) rectangle (1.6,0.35);
    \draw[fill=red!15] (2.2,-0.5) rectangle (2.55,0.9);
    \node[font=\tiny] at (0.8,0.17) {row $i$};
    \node[font=\tiny,rotate=90] at (2.37,0.2) {col $j$};
    \node[font=\tiny] at (1.9,0.17) {$\cdot=$};
    \node[draw,fill=violet!15,minimum size=0.35cm,font=\tiny] at (3.3,0.17) {$c_{ij}$};
  \end{scope}
\end{tikzpicture}
\caption{Matrix product: inner dimensions agree; entry $c_{ij}$ is the dot product of row $i$ of $A$ with column $j$ of $B$.}
\label{fig:matmul}
\end{figure}

Crucial facts: $AB\neq BA$ in general; $A(BC)=(AB)C$ (associative); $A(B+C)=AB+AC$; there is an identity $I_n$ with $I_nA=A$, $AI_n=A$.

\begin{example}[Composition viewpoint]
If $B:\mathbb{R}^p\to\mathbb{R}^n$ and $A:\mathbb{R}^n\to\mathbb{R}^m$ are linear maps, then $AB$ is the composition ``do $B$, then $A$.'' Order matters --- just as putting on socks then shoes $\neq$ shoes then socks. That is why $AB\neq BA$.
\end{example}

\subsection{Four Ways to See $A\mathbf{x}$}
For $A\in\mathbb{R}^{m\times n}$, $\mathbf{x}\in\mathbb{R}^n$:
\begin{enumerate}[label=(\roman*)]
    \item \textbf{Row dot:} $(A\mathbf{x})_i=\text{row}_i(A)\cdot\mathbf{x}$.
    \item \textbf{Column combination:} $A\mathbf{x}=x_1\mathbf{a}_1+\cdots+x_n\mathbf{a}_n$ where $\mathbf{a}_j$ is column $j$ of $A$. So $A\mathbf{x}$ lives in the span of the columns.
    \item \textbf{Linear map:} $\mathbf{x}\mapsto A\mathbf{x}$ is linear: $A(c\mathbf{x}+d\mathbf{y})=cA\mathbf{x}+dA\mathbf{y}$.
    \item \textbf{Outer sum:} $AB=\sum_{k}\mathbf{a}_k\mathbf{b}_k^T$ (sum of column-times-row rank-one pieces).
\end{enumerate}
View (ii) is the most useful: the columns tell you where the basis lands.

\begin{figure}[h]
\centering
\begin{tikzpicture}[scale=0.82,every node/.style={font=\scriptsize}]
  \draw[->,thick] (-0.3,0)--(3.2,0) node[right] {$x_1$};
  \draw[->,thick] (0,-0.3)--(0,3.0) node[above] {$x_2$};
  \draw[->,very thick,blue!65!black] (0,0)--(1.5,0.4) node[midway,below] {$\mathbf{a}_1$};
  \draw[->,very thick,red!65!black] (0,0)--(0.5,1.8) node[midway,left] {$\mathbf{a}_2$};
  \draw[->,very thick,violet!70!black] (0,0)--(2.0,2.2) node[midway,above right] {$A\mathbf{x}=x_1\mathbf{a}_1+x_2\mathbf{a}_2$};
  \draw[dashed,gray] (1.5,0.4)--(2.0,2.2)--(0.5,1.8);
  \fill[blue!10,opacity=0.5] (0,0)--(1.5,0.4)--(2.0,2.2)--(0.5,1.8)--cycle;
  \node[font=\tiny,fill=white] at (1.0,1.1) {parallelogram};
  \node[font=\tiny,fill=white] at (2.4,0.15) {$x_1=1,x_2=1$};
\end{tikzpicture}
\caption{$A\mathbf{x}$ as a linear combination of the columns $\mathbf{a}_1,\mathbf{a}_2$ of $A$. The set of all $A\mathbf{x}$ is the column span.}
\label{fig:col-combo}
\end{figure}

\section{Transpose and Special Matrices}
\label{sec:transpose}
\begin{definition}[Transpose]
$(A^T)_{ij}=A_{ji}$ --- flip across the diagonal. $(AB)^T=B^TA^T$, $(A^T)^T=A$.
\end{definition}
$A$ is \emph{symmetric} if $A^T=A$ (square, $a_{ij}=a_{ji}$); \emph{diagonal} if $a_{ij}=0$ for $i\neq j$; \emph{upper triangular} if $a_{ij}=0$ for $i>j$; \emph{identity} $I_n=\operatorname{diag}(1,\dots,1)$.

Trace: $\operatorname{tr}(A)=\sum_i a_{ii}$; $\operatorname{tr}(AB)=\operatorname{tr}(BA)$.

\begin{example}[Outer product]
For $\mathbf{u}\in\mathbb{R}^m$, $\mathbf{v}\in\mathbb{R}^n$, the outer product $\mathbf{u}\mathbf{v}^T\in\mathbb{R}^{m\times n}$ has rank one. Inner product $\mathbf{u}^T\mathbf{v}$ is $1\times1$ (a scalar). This distinction matters in ML: $\mathbf{x}\mathbf{x}^T$ is a covariance building block.
\end{example}

\begin{example}[Python --- matrix arithmetic]
\begin{lstlisting}[language=Python]
import numpy as np
A = np.array([[1,2,3],[4,5,6]], dtype=float)  # 2x3
B = np.array([[1,0],[0,1],[1,1]], dtype=float)  # 3x2
C = A @ B            # 2x2 matrix product
print(C)             # [[4,5],[10,11]]
print(A.T.shape)     # (3,2) -- transpose flips shape
# Column view: A @ x = x[0]*col0 + x[1]*col1 + x[2]*col2
x = np.array([1, -1, 2])
print(A @ x, 1*A[:,0] + -1*A[:,1] + 2*A[:,2])
\end{lstlisting}
\end{example}

% ============================================================

% --- Additional depth: Inverses and linear maps ---
\section{Matrix Inverses}
\label{sec:inverses}
$A\in\mathbb{R}^{n\times n}$ is \emph{invertible} (nonsingular) if there exists $A^{-1}$ with $AA^{-1}=A^{-1}A=I_n$. Then $A\mathbf{x}=\mathbf{b}$ has unique solution $\mathbf{x}=A^{-1}\mathbf{b}$, and $A$ represents a bijective linear map. Not all square matrices are invertible: $\begin{bmatrix}1&2\\2&4\end{bmatrix}$ squashes $\mathbb{R}^2$ onto a line.

\begin{theorem}[Invertibility checklist]
For $A\in\mathbb{R}^{n\times n}$, TFAE: (i) $A$ invertible; (ii) $\mathcal{N}(A)=\{\mathbf{0}\}$; (iii) $\mathcal{C}(A)=\mathbb{R}^n$; (iv) $\operatorname{rank}A=n$; (v) $\det A\neq0$ (Ch.~7); (vi) columns (equivalently rows) are independent.
\end{theorem}

$2\times2$ inverse: $\begin{bmatrix}a&b\\c&d\end{bmatrix}^{-1}=\frac1{ad-bc}\begin{bmatrix}d&-b\\-c&a\end{bmatrix}$ when $ad-bc\neq0$. For larger $n$, elimination on $[A\mid I]$ yields $[I\mid A^{-1}]$.

\section{Linear Maps and Matrix Representation}
\label{sec:linmap-detail}
A function $T:\mathbb{R}^n\to\mathbb{R}^m$ is \emph{linear} if $T(c\mathbf{u}+d\mathbf{v})=cT(\mathbf{u})+dT(\mathbf{v})$. Every such $T$ is $\mathbf{x}\mapsto A\mathbf{x}$ for a unique $A$ whose columns are $T(\mathbf{e}_j)$ ($\mathbf{e}_j$ standard basis). Composition $S\circ T$ corresponds to product $BA$.

\begin{example}[Rotation in $\mathbb{R}^2$]
Rotation by $\theta$ counter-clockwise: $R_\theta=\begin{bmatrix}\cos\theta&-\sin\theta\\\sin\theta&\cos\theta\end{bmatrix}$. Then $R_\theta R_\phi=R_{\theta+\phi}$ and $R_\theta^T=R_{-\theta}=R_\theta^{-1}$ (orthogonal matrix).
\end{example}



\begin{remark}[Why $AB\neq BA$ matters in code]
In NumPy, \texttt{A @ B} and \texttt{B @ A} may even have different shapes. Computationally, $(AB)\mathbf{x}=A(B\mathbf{x})$ suggests parenthesising to minimise flops: for $A\in\mathbb{R}^{10\times100}$, $B\in\mathbb{R}^{100\times10}$, $\mathbf{x}\in\mathbb{R}^{10}$, compute $B\mathbf{x}$ first ($100\times10$ by $10$ $\to O(1000)$) then $A(\cdot)$ rather than forming $AB$ ($10\times10$ costs $O(10\cdot100\cdot10)=10^4$). Order and bracketing affect both meaning and cost.
\end{remark}

\begin{example}[Sparse matrices]
Many computing matrices are sparse (mostly zeros): adjacency of a graph, finite differences. Storing only nonzeros and multiplying by iterating over nonzeros saves $O(n^2)$ to $O(\text{nnz})$. \texttt{scipy.sparse} does this; dense \texttt{@} would waste time on zeros.
\end{example}



\begin{example}[Hadamard vs.\ matrix product]
Entrywise (Hadamard) $A\circ B=[a_{ij}b_{ij}]$ is sometimes written $A*B$ in code (e.g.\ NumPy \texttt{A*B} for arrays). It is commutative and $O(mn)$, but it is \emph{not} the matrix product $AB$. Mixing them is a common bug: check whether you need \texttt{A @ B} (composition) or \texttt{A * B} (masking/scaling).
\end{example}

\begin{remark}[Transpose and symmetry in computing]
Storing $A^T$ is often free (stride trick, not a copy) in NumPy; \texttt{A.T} is a view. Symmetry $A^T=A$ halves storage; triangular storage saves half the entries. Numerical libraries exploit this: Cholesky for symmetric positive-definite $A$ is $2\times$ faster than $LU$.
\end{remark}


\section{Chapter Notes}
Matrix multiplication as composition explains non-commutativity and associativity at once. The column picture anticipates column space and rank (Ch.~5) and the $A=LU$ factorisation (Ch.~3). For cost, $m\times n$ times $n\times p$ is $O(mnp)$ naively; Strassen (Ch.~3 notes) does better.

% ============================================================
\section{Exercises}
\label{sec:ch02-exercises}
\begin{enumerate}[label=E2.\arabic*, leftmargin=*, itemsep=0.8em]
\item \textbf{Product by hand.} Let $A=\begin{bmatrix}2&-1&0\\1&3&2\end{bmatrix}$, $B=\begin{bmatrix}1&2\\0&-1\\3&1\end{bmatrix}$. Compute $AB$ and $BA$ (if defined). Verify $(AB)^T=B^TA^T$ numerically.
\item \textbf{Column view.} Write $A=\begin{bmatrix}1&2\\3&4\\5&6\end{bmatrix}$ and $\mathbf{x}=[2,-1]^T$. Compute $A\mathbf{x}$ both as row-dots and as a combination of columns. What is the set $\{A\mathbf{x}:\mathbf{x}\in\mathbb{R}^2\}$ geometrically?
\item \textbf{Algebra.} Prove: if $A,B$ symmetric and $AB=BA$ then $AB$ is symmetric. Give a $2\times2$ counterexample when they do not commute.
\item \textbf{Block multiply.} Let $M=\begin{bmatrix}I&A\\0&I\end{bmatrix}$ where $A$ is $n\times n$. Find $M^2$ and $M^{-1}$ in block form and verify by block multiplication.
\item \textbf{Coding.} Benchmark \texttt{A @ B} for $n\times n$ random matrices ($n=100,500,1000$) with NumPy. Plot time vs.\ $n$ on a log--log scale and estimate the empirical exponent. Why is NumPy far faster than triple loops?
\end{enumerate}
% End of Chapter 2

% SC1004 Chapter 3: Linear Systems & Gauss Elimination (0->100 ground-up)
\chapter{Linear Systems and Gaussian Elimination}
\label{ch:systems}

\begin{quote}
\emph{``Elimination is just organised subtraction.''} A system of equations is a puzzle of balances; Gaussian elimination solves it by the same moves you would use on paper --- swap, scale, subtract --- written as row operations on a matrix.
\end{quote}

\noindent\fbox{\parbox{\dimexpr\linewidth-2\fboxsep-2\fboxrule}{%
\textbf{From zero: no systems background needed.} We start with two lines in the plane, see the three things that can happen (one point, no point, a whole line), then grow to $m\times n$ and build an algorithm that handles them all.}}
\vspace{0.4em}

\section*{Learning Outcomes}
After this chapter you will be able to:
\begin{enumerate}[label=(L\arabic*)]
    \item write a linear system as $A\mathbf{x}=\mathbf{b}$ and as an augmented matrix $[A\mid\mathbf{b}]$;
    \item carry out Gaussian elimination to row-echelon and Gauss--Jordan forms;
    \item read off existence, uniqueness, and parametric solution sets;
    \item relate pivots, free variables, and rank to the geometry;
    \item factor $A=LU$ and use it to solve efficiently.
\end{enumerate}

% ============================================================
\section{Systems as Matrices}
\label{sec:systems-matrix}
A \emph{linear equation} is $a_1x_1+\cdots+a_nx_n=b$. A \emph{system} of $m$ such equations is
\[
  A\mathbf{x}=\mathbf{b},\quad A\in\mathbb{R}^{m\times n},\ \mathbf{x}\in\mathbb{R}^n,\ \mathbf{b}\in\mathbb{R}^m.
\]
The \emph{augmented matrix} $[A\mid\mathbf{b}]$ keeps all the data in one table.

\begin{figure}[h]
\centering
\begin{tikzpicture}[scale=0.85,every node/.style={font=\scriptsize}]
  \draw[->,thick] (-0.4,0)--(3.2,0) node[right] {$x$};
  \draw[->,thick] (0,-0.4)--(0,3.0) node[above] {$y$};
  % one intersection
  \begin{scope}
    \draw[very thick,blue!60!black] (-0.2,0.3)--(2.6,2.6) node[right] {$L_1$};
    \draw[very thick,red!60!black] (-0.2,2.4)--(2.6,0.5) node[right] {$L_2$};
    \fill[green!60!black] (1.35,1.45) circle (2.5pt) node[above,font=\tiny] {unique};
    \node[font=\footnotesize\bfseries] at (1.3,-0.7) {One solution};
  \end{scope}
  \begin{scope}[xshift=4.0cm]
    \draw[very thick,blue!60!black] (-0.2,0.3)--(2.6,2.6);
    \draw[very thick,red!60!black] (-0.2,0.9)--(2.6,3.2);
    \node[font=\tiny,fill=white] at (1.2,1.5) {parallel};
    \node[font=\footnotesize\bfseries] at (1.3,-0.7) {No solution};
  \end{scope}
  \begin{scope}[xshift=8.0cm]
    \draw[very thick,blue!60!black] (-0.2,0.3)--(2.6,2.6);
    \draw[very thick,red!60!black,dashed] (-0.2,0.3)--(2.6,2.6);
    \node[font=\tiny,fill=white] at (1.2,1.2) {coincident};
    \node[font=\footnotesize\bfseries] at (1.3,-0.7) {Infinitely many};
  \end{scope}
\end{tikzpicture}
\caption{Three geometries for two equations in two unknowns. With more variables the same trichotomy holds.}
\label{fig:three-cases}
\end{figure}

For $n>2$ we cannot draw, but algebra does not care: the same three outcomes occur.

\section{Row Operations and Echelon Form}
\label{sec:row-ops}
Elimination uses three \emph{elementary row operations} on $[A\mid\mathbf{b}]$, each reversible and preserving the solution set:
\begin{enumerate}[label=(R\arabic*)]
    \item Swap two rows: $R_i\leftrightarrow R_j$.
    \item Scale a row: $R_i\gets cR_i$, $c\neq0$.
    \item Add a multiple: $R_i\gets R_i+cR_j$, $i\neq j$.
\end{enumerate}

\begin{definition}[Row-echelon form (REF)]
A matrix is in REF if (i) all zero rows are at the bottom, (ii) each leading entry (pivot) is to the right of the one above, (iii) pivots may be any nonzero. In \emph{reduced} REF (RREF), each pivot is $1$ and is the only nonzero in its column.
\end{definition}

\begin{figure}[h]
\centering
\begin{tikzpicture}[scale=0.78,every node/.style={font=\scriptsize}]
  \foreach \r in {0,...,2} \foreach \c in {0,...,4} {
    \pgfmathsetmacro{\shade}{(\r==0 && \c==0) || (\r==1 && \c==1) || (\r==2 && \c==3) ? 1 : 0}
    \ifnum\shade=1
      \node[draw,thick,minimum width=0.78cm,minimum height=0.55cm,fill=orange!35] at (\c*0.88,\r*-0.65) {$\star$};
    \else
      \pgfmathparse{(\r==0 || (\r==1 && \c>=1) || (\r==2 && \c>=3)) ? 1 : 0}
      \ifnum\pgfmathresult=1
        \node[draw,minimum width=0.78cm,minimum height=0.55cm,fill=blue!7] at (\c*0.88,\r*-0.65) {$\cdot$};
      \else
        \node[draw,minimum width=0.78cm,minimum height=0.55cm,fill=gray!15] at (\c*0.88,\r*-0.65) {$0$};
      \fi
    \fi
  }
  \draw[very thick,red!60!black] (-0.45,0.32)--(4.2,-1.55) node[right,font=\tiny] {staircase of pivots};
  \node[font=\tiny] at (2.0,-2.2) {REF: pivots $\star$ step right; zeros below staircase};
  \node[draw,rounded corners=2pt,fill=green!12,font=\tiny,align=center] at (5.6,-0.65) {pivots $=$\\leading entries};
\end{tikzpicture}
\caption{Shape of a row-echelon matrix. The staircase separates the pivot columns from free columns.}
\label{fig:echelon}
\end{figure}

\emph{Gaussian elimination:} work column by column, pick a pivot, clear below with (R3), move to next row/column. \emph{Gauss--Jordan} continues upward to clear above pivots, reaching RREF.

\begin{example}[Elimination by hand]
Solve
\[
\begin{cases}
  x+2y+z=4\\ 2x+5y+z=7\\ x+y- z=0
\end{cases}\quad
\left[\begin{array}{ccc|c}1&2&1&4\\2&5&1&7\\1&1&-1&0\end{array}\right]
\xrightarrow{R_2-2R_1,\;R_3-R_1}
\left[\begin{array}{ccc|c}1&2&1&4\\0&1&-1&-1\\0&-1&-2&-4\end{array}\right]
\xrightarrow{R_3+R_2}
\left[\begin{array}{ccc|c}1&2&1&4\\0&1&-1&-1\\0&0&-3&-5\end{array}\right].
\]
Back substitution: $z=5/3$, $y=-1+z=2/3$, $x=4-2y-z=1$. Unique solution $\mathbf{x}=[1,2/3,5/3]^T$. If a row became $[0\;0\;0\mid1]$, inconsistency --- no solution.
\end{example}

\section{Solution Sets: Pivots vs.\ Free Variables}
\label{sec:solution-sets}
Each column with a pivot is a \emph{pivot column}; the rest are \emph{free}. Free variables become parameters.

\begin{theorem}[Existence and uniqueness]
$A\mathbf{x}=\mathbf{b}$ has a solution iff $\operatorname{rank}(A)=\operatorname{rank}([A\mid\mathbf{b}])$ (no pivot in the augmented column). It is unique iff every variable column is a pivot column (no free variables).
\end{theorem}

When free variables exist, the solution is an affine subspace:
\[
  \mathbf{x}=\mathbf{p}+t_1\mathbf{v}_1+\cdots+t_k\mathbf{v}_k,
\]
$\mathbf{p}$ a particular solution, $\{\mathbf{v}_i\}$ span the homogeneous solutions $A\mathbf{x}=\mathbf{0}$ (the null space, Ch.~4).

\begin{figure}[h]
\centering
\begin{tikzpicture}[scale=0.80,every node/.style={font=\scriptsize}]
  \draw[->,thick] (-1.2,0)--(3.5,0) node[right] {$x$};
  \draw[->,thick] (0,-1.2)--(0,3.0) node[above] {$y$};
  \draw[->,thick] (-0.8,-0.8)--(1.5,2.2) node[above] {$z$};
  \fill[blue!12] (-0.8,-0.3)--(3.2,0.2)--(2.0,2.6)--(-1.2,2.1)--cycle;
  \draw[very thick,blue!60!black] (-0.8,-0.3)--(3.2,0.2)--(2.0,2.6)--(-1.2,2.1)--cycle;
  \fill[red!70!black] (0.8,0.9) circle (2.5pt) node[above right] {$\mathbf{p}$};
  \draw[->,very thick,orange!75!black] (0.8,0.9)--(1.9,1.1) node[midway,above] {$\mathbf{v}_1$};
  \draw[->,very thick,violet!70!black] (0.8,0.9)--(0.4,1.8) node[midway,left] {$\mathbf{v}_2$};
  \node[font=\tiny,fill=white] at (1.0,-0.6) {Plane of solutions when 1 pivot missing in $\mathbb{R}^3$};
\end{tikzpicture}
\caption{With free variables, solutions form $\mathbf{p}+\operatorname{span}\{\mathbf{v}_i\}$ --- a point, line, plane, or hyperplane.}
\label{fig:solution-affine}
\end{figure}

\section{LU Factorisation}
\label{sec:lu}
Elimination records its multipliers in $L$ (unit lower-triangular) so that $A=LU$ with $U$ the REF. Then $A\mathbf{x}=\mathbf{b}$ becomes $L\mathbf{y}=\mathbf{b}$ (forward substitution) then $U\mathbf{x}=\mathbf{y}$ (back substitution) --- $O(n^2)$ per new $\mathbf{b}$ after one $O(n^3)$ factorisation. With row swaps, $PA=LU$.

\begin{example}[Python --- elimination and LU]
\begin{lstlisting}[language=Python]
import numpy as np
from scipy.linalg import lu
A = np.array([[1.,2.,1.],[2.,5.,1.],[1.,1.,-1.]])
b = np.array([4.,7.,0.])
x = np.linalg.solve(A, b)
print(x)  # [1. 0.666... 1.666...]
P, L, U = lu(A)
print(np.allclose(P@L@U, A))  # True
\end{lstlisting}
If \texttt{np.linalg.matrix\_rank(A)} $<n$, expect either no or infinitely many solutions --- check augmented rank.
\end{example}

% ============================================================

% --- Additional depth: Rank, nullity preview, and iterative view ---
\section{Rank and Its Geometry}
\label{sec:rank-geom}
The number of pivots $r=\operatorname{rank}(A)$ counts independent equations (or independent columns). Full row rank ($r=m$) means $A\mathbf{x}=\mathbf{b}$ solvable for every $\mathbf{b}$; full column rank ($r=n$) means at most one solution for each $\mathbf{b}$. When $r<\min(m,n)$, both existence and uniqueness can fail depending on $\mathbf{b}$.

\begin{theorem}[Rank of products and sums]
$\operatorname{rank}(AB)\le\min(\operatorname{rank}A,\operatorname{rank}B)$; $\operatorname{rank}(A+B)\le\operatorname{rank}A+\operatorname{rank}B$. Row operations preserve rank (they are left-multiplication by invertible $E$).
\end{theorem}

\section{Computational Cost and Pivoting}
\label{sec:pivot-cost}
Elimination on $n\times n$ costs $\sim\frac23n^3$ flops (floating-point ops); forward/back substitution $O(n^2)$. Partial pivoting (swap to bring largest $|a_{ik}|$ to pivot) avoids division by tiny numbers and controls error growth (factor $\le2^{n-1}$ worst-case, tiny in practice). Complete pivoting (row+column swaps) is more stable but costlier. Libraries (LAPACK) use blocked $LU$ for cache efficiency.

\begin{example}[When pivoting matters]
$A=\begin{bmatrix}10^{-12}&1\\1&1\end{bmatrix}$ without pivoting: first pivot $10^{-12}$ causes catastrophic cancellation in $R_2\gets R_2-10^{12}R_1$. Swapping rows first gives pivot $1$, stable elimination.
\end{example}



\begin{remark}[Geometric reading of REF]
Each pivot column corresponds to a direction that is constrained; each free column is a degree of freedom. So $\operatorname{rank}A$ is the dimension of the column space (how many independent constraints), and $n-\operatorname{rank}A$ is the dimension of the solution line/plane when $\mathbf{b}$ is consistent. This foreshadows Ch.~5's rank--nullity.
\end{remark}

\begin{example}[No solution vs.\ infinitely many]
$A=\begin{bmatrix}1&1\\2&2\end{bmatrix}$, $\mathbf{b}_1=[1,3]^T$: augmented REF $\begin{bmatrix}1&1&1\\0&0&1\end{bmatrix}$ --- pivot in augmented column, no solution (parallel lines). $\mathbf{b}_2=[1,2]^T$: $\begin{bmatrix}1&1&1\\0&0&0\end{bmatrix}$ --- free $x_2=t$, solutions $[1-t,t]^T$ (coincident lines).
\end{example}


\section{Chapter Notes}
Gauss described elimination circa 1809 (least squares); Jordan pushed to RREF. $LU$ is how libraries actually solve $A\mathbf{x}=\mathbf{b}$; QR (Ch.~6) is the stable alternative for least squares. Pivoting (partial: largest in column) controls round-off.

% ============================================================
\section{Exercises}
\label{sec:ch03-exercises}
\begin{enumerate}[label=E3.\arabic*, leftmargin=*, itemsep=0.8em]
\item \textbf{REF by hand.} Reduce $\begin{bmatrix}1&2&3&9\\2&4&5&13\\1&0&-1&-1\end{bmatrix}$ to REF and RREF. Identify pivots and free variables; write all solutions.
\item \textbf{Trichotomy.} Give matrices $A$ and vectors $\mathbf{b}_1,\mathbf{b}_2,\mathbf{b}_3$ ($2\times2$) exhibiting the three cases: unique solution, no solution, infinitely many. Explain via ranks of $A$ and $[A\mid\mathbf{b}]$.
\item \textbf{Homogeneous.} Show the solution set of $A\mathbf{x}=\mathbf{0}$ is closed under addition and scaling. Find a basis for it when $A=\begin{bmatrix}1&2&-1&0\\0&1&1&1\end{bmatrix}$.
\item \textbf{LU.} Compute $L$ and $U$ (no pivoting) for $A=\begin{bmatrix}2&1&1\\4&5&2\\2&1&-1\end{bmatrix}$ by recording multipliers. Verify $A=LU$ and solve $A\mathbf{x}=[1,2,0]^T$ via $L,U$.
\item \textbf{Coding.} Write \texttt{rref(A)} (Gauss--Jordan, partial pivoting, tolerance $10^{-10}$) returning RREF and pivot columns. Test on random $3\times5$ matrices and compare pivot columns with \texttt{np.linalg.matrix\_rank}.
\end{enumerate}
% End of Chapter 3

% SC1004 Chapter 4: Vector Spaces (0->100 ground-up)
\chapter{Vector Spaces}
\label{ch:vectorspaces}

\begin{quote}
\emph{``A vector space is any place where arrows still make sense.''} We abstract $\mathbb{R}^n$ to uncover spaces of matrices, polynomials, signals, and images --- all governed by the same eight rules.
\end{quote}

\noindent\fbox{\parbox{\dimexpr\linewidth-2\fboxsep-2\fboxrule}{%
\textbf{From zero: axioms after examples.} We have lived in $\mathbb{R}^n$ for three chapters. Now we ask what makes it tick, write the rules down, and test which other collections obey them. Every axiom is motivated by a picture or a failure.}}
\vspace{0.4em}

\section*{Learning Outcomes}
After this chapter you will be able to:
\begin{enumerate}[label=(L\arabic*)]
    \item state the eight vector-space axioms and test a set against them;
    \item recognise subspaces and verify them via the subspace test;
    \item describe column space, null space, row space, and their relations;
    \item decide membership in a span and express vectors as combinations;
    \item connect subspaces to solutions of $A\mathbf{x}=\mathbf{b}$ and $A\mathbf{x}=\mathbf{0}$.
\end{enumerate}

% ============================================================
\section{Why Abstract?}
\label{sec:why-abstract}
In $\mathbb{R}^n$ we add arrows tip-to-tail and stretch them. But we also add \emph{polynomials} $(p+q)(x)=p(x)+q(x)$, \emph{matrices}, \emph{audio signals} (add waveforms, scale volume), and \emph{RGB images}. Each supports the same two operations. Rather than prove everything four times, we isolate the common rules.

\section{The Axioms}
\label{sec:axioms}
A \emph{vector space} over $\mathbb{R}$ is a set $V$ with addition $+$ and scalar multiplication $\cdot$ satisfying, for all $\mathbf{u},\mathbf{v},\mathbf{w}\in V$, $c,d\in\mathbb{R}$:
\begin{enumerate}[label=(V\arabic*)]
    \item $\mathbf{u}+\mathbf{v}\in V$ (closure under $+$),\quad $c\mathbf{u}\in V$ (closure under scaling);
    \item $\mathbf{u}+\mathbf{v}=\mathbf{v}+\mathbf{u}$; $(\mathbf{u}+\mathbf{v})+\mathbf{w}=\mathbf{u}+(\mathbf{v}+\mathbf{w})$;
    \item there exists $\mathbf{0}\in V$ with $\mathbf{v}+\mathbf{0}=\mathbf{v}$;
    \item for each $\mathbf{v}$ there is $-\mathbf{v}$ with $\mathbf{v}+(-\mathbf{v})=\mathbf{0}$;
    \item $c(\mathbf{u}+\mathbf{v})=c\mathbf{u}+c\mathbf{v}$;\; $(c+d)\mathbf{v}=c\mathbf{v}+d\mathbf{v}$;
    \item $c(d\mathbf{v})=(cd)\mathbf{v}$;\; $1\mathbf{v}=\mathbf{v}$.
\end{enumerate}
Elements of $V$ are called \emph{vectors} even when they look like matrices or functions.

\begin{figure}[h]
\centering
\begin{tikzpicture}[scale=0.82,every node/.style={font=\scriptsize}]
  \node[draw,thick,rounded corners=4pt,fill=blue!8,minimum width=3.2cm,minimum height=2.0cm] (V) at (0,0) {};
  \node[font=\footnotesize\bfseries] at (0,1.25) {Vector space $V$};
  \fill[blue!55!black] (-0.5,0.3) circle (2pt) node[above] {$\mathbf{u}$};
  \fill[red!65!black] (0.7,0.2) circle (2pt) node[above] {$\mathbf{v}$};
  \fill[violet!65!black] (0.15,0.75) circle (2pt) node[above] {$\mathbf{u}+\mathbf{v}$};
  \draw[->,thick,blue!50!black,dashed] (-0.5,0.3)--(0.15,0.75);
  \draw[->,thick,red!50!black,dashed] (0.7,0.2)--(0.15,0.75);
  \fill[orange!75!black] (0.4,-0.55) circle (2pt) node[below] {$c\mathbf{u}$};
  \draw[->,thick,orange!60!black] (0,0)--(0.4,-0.55);
  \node[font=\tiny,fill=white] at (0,-1.25) {closed under $+$ and scaling};

  \node[draw,thick,rounded corners=4pt,fill=green!8,minimum width=2.6cm,minimum height=2.0cm] at (4.5,0) {};
  \node[font=\footnotesize\bfseries] at (4.5,1.25) {Examples};
  \node[font=\tiny,align=left] at (4.5,0.5) {$\mathbb{R}^n$ --- arrows};
  \node[font=\tiny,align=left] at (4.5,0.1) {$\mathbb{R}^{m\times n}$ --- matrices};
  \node[font=\tiny,align=left] at (4.5,-0.3) {$\mathcal{P}_d$ --- polynomials};
  \node[font=\tiny,align=left] at (4.5,-0.7) {$C[0,1]$ --- signals};
  \node[font=\tiny,align=left] at (4.5,-1.1) {$\{\mathbf{0}\}$ --- trivial};
\end{tikzpicture}
\caption{Left: closure means sums and scalings never leave $V$. Right: diverse examples obeying the same axioms.}
\label{fig:vecspace-closure}
\end{figure}

\begin{example}[Polynomials]
$\mathcal{P}_2=\{a_0+a_1x+a_2x^2:a_i\in\mathbb{R}\}$ is a vector space of dimension $3$ (basis $1,x,x^2$). The set $\{p:p(0)=1\}$ is \emph{not} a subspace (fails $\mathbf{0}$ and closure --- sum has $p(0)=2$).
\end{example}

\begin{example}[Non-example]
The first quadrant $\{(x,y):x\ge0,y\ge0\}\subset\mathbb{R}^2$ is not a vector space: $(1,1)$ is in it but $-(1,1)=(-1,-1)$ is not (fails (V4) and scaling by $-1$).
\end{example}

\section{Subspaces}
\label{sec:subspaces}
\begin{definition}[Subspace]
$W\subseteq V$ is a \emph{subspace} if it is itself a vector space under the same $+$ and $\cdot$. Equivalently (subspace test):
\[
  W\neq\emptyset,\quad \mathbf{u},\mathbf{v}\in W\Rightarrow \mathbf{u}+\mathbf{v}\in W,\quad c\in\mathbb{R},\mathbf{u}\in W\Rightarrow c\mathbf{u}\in W.
\]
In particular $\mathbf{0}\in W$ always.
\end{definition}

\begin{figure}[h]
\centering
\begin{tikzpicture}[scale=0.82,every node/.style={font=\scriptsize}]
  \draw[->,thick] (-0.4,0)--(3.2,0) node[right] {$x$};
  \draw[->,thick] (0,-0.4)--(0,3.0) node[above] {$y$};
  \draw[->,thick] (-0.6,-0.6)--(1.8,2.2) node[above] {$z$};
  % plane through origin = subspace
  \fill[blue!12] (-0.9,-0.2)--(2.8,0.3)--(1.6,2.5)--(-1.0,2.0)--cycle;
  \draw[very thick,blue!60!black] (-0.9,-0.2)--(2.8,0.3)--(1.6,2.5)--(-1.0,2.0)--cycle;
  \fill[red!70!black] (0,0) circle (2.5pt) node[below left] {$\mathbf{0}$};
  \draw[->,very thick,orange!75!black] (0,0)--(1.2,0.55) node[midway,below] {$\mathbf{u}$};
  \draw[->,very thick,violet!70!black] (0,0)--(0.2,1.1) node[midway,left] {$\mathbf{v}$};
  \draw[->,thick,green!60!black,dashed] (0,0)--(1.4,1.65) node[midway,above right] {$\mathbf{u}+\mathbf{v}$ still in plane};
  \node[font=\tiny,fill=white] at (1.0,-0.7) {Plane through $\mathbf{0}$: subspace of $\mathbb{R}^3$};

  \begin{scope}[xshift=5.5cm]
    \fill[red!12] (-0.9,0.8)--(2.8,1.3)--(1.6,3.5)--(-1.0,3.0)--cycle;
    \draw[very thick,red!60!black] (-0.9,0.8)--(2.8,1.3)--(1.6,3.5)--(-1.0,3.0)--cycle;
    \fill[gray!60!black] (0,0) circle (2pt) node[below] {$\mathbf{0}\notin$ plane};
    \node[font=\tiny,fill=white] at (0.8,0.3) {shifted plane};
    \node[font=\tiny,align=center] at (0.9,-0.6) {Plane not through $\mathbf{0}$:\\not a subspace};
  \end{scope}
\end{tikzpicture}
\caption{Left: a subspace must contain $\mathbf{0}$ and be closed (plane through origin). Right: a translate fails the test.}
\label{fig:subspace-plane}
\end{figure}

\section{Four Fundamental Subspaces of a Matrix}
\label{sec:four-subspaces}
For $A\in\mathbb{R}^{m\times n}$:
\begin{itemize}
    \item \textbf{Column space} $\mathcal{C}(A)=\{A\mathbf{x}:\mathbf{x}\in\mathbb{R}^n\}\subseteq\mathbb{R}^m$ --- all $\mathbf{b}$ for which $A\mathbf{x}=\mathbf{b}$ is solvable. Span of the columns.
    \item \textbf{Null space} $\mathcal{N}(A)=\{\mathbf{x}:A\mathbf{x}=\mathbf{0}\}\subseteq\mathbb{R}^n$ --- homogeneous solutions (Ch.~3). Always a subspace.
    \item \textbf{Row space} $\mathcal{C}(A^T)\subseteq\mathbb{R}^n$ --- span of the rows.
    \item \textbf{Left null space} $\mathcal{N}(A^T)\subseteq\mathbb{R}^m$.
\end{itemize}
$\mathcal{C}(A)$ and $\mathcal{N}(A^T)$ live in $\mathbb{R}^m$; $\mathcal{N}(A)$ and $\mathcal{C}(A^T)$ in $\mathbb{R}^n$. Row operations preserve $\mathcal{N}(A)$ and row space but change $\mathcal{C}(A)$ --- a common pitfall.

\begin{example}[Membership]
Is $\mathbf{b}=[1,2,3]^T$ in $\mathcal{C}(A)$ for $A=\begin{bmatrix}1&2\\2&4\\3&6\end{bmatrix}$? Second column is $2\times$ first, so $\mathcal{C}(A)=\operatorname{span}\{[1,2,3]^T\}$. Since $\mathbf{b}=1\cdot[1,2,3]^T$, yes. $\mathbf{b}'=[1,0,0]^T$ is not.
\end{example}

\begin{example}[Python --- null space]
\begin{lstlisting}[language=Python]
import numpy as np
from scipy.linalg import null_space
A = np.array([[1.,2.,1.],[2.,4.,2.],[1.,1.,0.]])
print("rank", np.linalg.matrix_rank(A))
NS = null_space(A)   # orthonormal basis for N(A)
print(NS.T)          # each row is a basis vector
print(A @ NS)        # ~ 0 (up to round-off)
\end{lstlisting}
\end{example}

% ============================================================

% --- Additional depth: Span, sum, and direct sum ---
\section{Span and Spanning Sets}
\label{sec:span-detail}
For $S=\{\mathbf{v}_1,\dots,\mathbf{v}_k\}\subseteq V$, $\operatorname{span}(S)=\{\sum c_i\mathbf{v}_i:c_i\in\mathbb{R}\}$ is the smallest subspace containing $S$ (intersection of all subspaces containing $S$). $S$ \emph{spans} $V$ if $\operatorname{span}(S)=V$. Spanning is the opposite of independence: too many vectors guarantees dependence, too few cannot span.

\begin{theorem}[Spanning test via matrices]
Fix a basis to coordinatize $V\cong\mathbb{R}^n$. Then $S$ spans $V$ iff the matrix with columns $[{\mathbf{v}_i}]$ has rank $n$ (every $\mathbf{b}$ is a combination).
\end{theorem}

\section{Sums and Direct Sums}
\label{sec:sums}
For subspaces $U,W\subseteq V$, $U+W=\{\mathbf{u}+\mathbf{w}:\mathbf{u}\in U,\mathbf{w}\in W\}$ is a subspace; $U\cap W$ is too. The sum is \emph{direct}, $V=U\oplus W$, if $V=U+W$ and $U\cap W=\{\mathbf{0}\}$ --- then every $\mathbf{v}\in V$ has a unique $\mathbf{u}+\mathbf{w}$ decomposition. Example: $\mathbb{R}^3=\text{$xy$-plane}\oplus\text{$z$-axis}$.

\begin{example}[Polynomial split]
$\mathcal{P}_3=\{p:p(0)=0\}\oplus\{ \text{constants}\}$: every $p(x)=a_0+a_1x+a_2x^2+a_3x^3$ uniquely as $(a_1x+a_2x^2+a_3x^3)+a_0$. Dimensions add: $3+1=4=\dim\mathcal{P}_3$.
\end{example}

\section{Coordinates Revisited}
\label{sec:coords-vec}
Fix a basis $B=\{\mathbf{b}_1,\dots,\mathbf{b}_n\}$ of $V$. The map $\mathbf{v}\mapsto[\mathbf{v}]_B\in\mathbb{R}^n$ (coordinates) is a linear isomorphism: it preserves $+$ and scaling and is bijective. So every $n$-dimensional real vector space is ``just'' $\mathbb{R}^n$ in disguise --- but the disguise matters for interpretation (e.g.\ polynomial vs.\ signal).



\begin{example}[Function space: audio signals]
Let $V=\{f:[0,1]\to\mathbb{R}\text{ continuous}\}$ with $(f+g)(t)=f(t)+g(t)$, $(cf)(t)=c\,f(t)$. This is infinite-dimensional: $\{1,t,t^2,\dots\}$ are independent but never span all continuous signals (need Fourier or polynomial approximations). Subspace $W=\{f:f(0)=0\}$ contains $\mathbf{0}$ and is closed; $W'=\{f:f(0)=1\}$ is not.
\end{example}

\begin{remark}[Why $\mathbf{0}\in W$ is non-negotiable]
If $W$ is closed under scaling, $0\cdot\mathbf{w}=\mathbf{0}\in W$ for any $\mathbf{w}\in W$ (provided $W\neq\emptyset$). So checking $\mathbf{0}\in W$ is the fastest way to kill a non-subspace: the shifted plane, the set $\det=1$, the unit circle all fail instantly.
\end{remark}



\section{Linear Combinations and Spanning Revisited}
\label{sec:lin-comb-detail}
A \emph{linear combination} of $\mathbf{v}_1,\dots,\mathbf{v}_k$ is $\sum c_i\mathbf{v}_i$. The set of all such is $\operatorname{span}\{\mathbf{v}_i\}$, already a subspace. To test if $\mathbf{b}\in\operatorname{span}\{\mathbf{v}_i\}$, solve $[\mathbf{v}_1\cdots\mathbf{v}_k]\mathbf{x}=\mathbf{b}$ --- this is exactly $A\mathbf{x}=\mathbf{b}$ from Ch.~3. Geometry (Ch.~1) and algebra (Ch.~2) meet: $\mathbf{b}$ is in the span iff it is a column combination, i.e.\ $\mathbf{b}\in\mathcal{C}(A)$.

\begin{example}[Polynomial coordinates]
In $\mathcal{P}_2$ with basis $\{1,x,x^2\}$, $p(x)=3-x+2x^2$ has $[p]_B=[3,-1,2]^T$. Addition and scaling act entrywise on coordinates, so $\mathcal{P}_2\cong\mathbb{R}^3$ as vector spaces --- but $p$ is a function, not an arrow, reminding us the isomorphism is interpretation-dependent.
\end{example}



\begin{remark}[Subspace vs.\ subset: a one-line test]
To kill a non-subspace, find one closure failure. The set $\{A:\det A=1\}$ fails scaling ($2A$ has $\det 2^n$); the unit circle fails addition ($(1,0)+(0,1)$ outside). One counterexample suffices.
\end{remark}


\section{Chapter Notes}
Abstract vector spaces (Peano, 1888; Banach, 1920s) let one theorem serve many domains. Subspaces are the ``lines/planes through the origin'' in any $V$. The four subspaces reappear in Ch.~6 as orthogonal complements.

% ============================================================
\section{Exercises}
\label{sec:ch04-exercises}
\begin{enumerate}[label=E4.\arabic*, leftmargin=*, itemsep=0.8em]
\item \textbf{Axiom check.} Which of these are vector spaces? (a) $\{(x,y):xy\ge0\}$ in $\mathbb{R}^2$, (b) $\{p\in\mathcal{P}_2:p(1)=0\}$, (c) $\{A\in\mathbb{R}^{2\times2}:\det A=0\}$. Justify via the subspace test or a counterexample.
\item \textbf{Span membership.} Let $A=\begin{bmatrix}1&0&2\\0&1&1\\1&1&3\end{bmatrix}$. Is $\mathbf{b}=[3,2,5]^T$ in $\mathcal{C}(A)$? Is $\mathbf{x}=[2,-1,1]^T$ in $\mathcal{N}(A)$? Decide by solving $A\mathbf{x}=\mathbf{b}$ / $A\mathbf{x}=\mathbf{0}$.
\item \textbf{Null space.} Find $\mathcal{N}(A)$ for $A=\begin{bmatrix}1&2&3\\2&4&6\end{bmatrix}$ explicitly as $\operatorname{span}\{\mathbf{v}_1,\mathbf{v}_2\}$. Verify closure under addition.
\item \textbf{Intersection.} Prove the intersection of two subspaces of $V$ is a subspace. Is the union always a subspace? Give a counterexample in $\mathbb{R}^2$.
\item \textbf{Coding.} Write \texttt{is\_in\_colspace(A,b,tol=1e-9)} using rank: $\mathbf{b}\in\mathcal{C}(A)$ iff $\operatorname{rank}(A)=\operatorname{rank}([A\mid\mathbf{b}])$. Test on random $4\times3$ matrices.
\end{enumerate}
% End of Chapter 4

% SC1004 Chapter 5: Independence, Basis, Dimension (0->100 ground-up)
\chapter{Linear Independence, Basis, and Dimension}
\label{ch:independence}

\begin{quote}
\emph{``Independence means no vector is a passenger --- each pulls in a new direction.''} This chapter asks when a set is redundant, how to extract the essentials, and how many essentials any space needs.
\end{quote}

\noindent\fbox{\parbox{\dimexpr\linewidth-2\fboxsep-2\fboxrule}{%
\textbf{From zero: redundancy before rank.} We build independence from the question ``can one vector be made from the others?'' then define basis as a minimal spanning set and dimension as its size --- proved well-defined.}}
\vspace{0.4em}

\section*{Learning Outcomes}
After this chapter you will be able to:
\begin{enumerate}[label=(L\arabic*)]
    \item test a set for linear independence via $A\mathbf{x}=\mathbf{0}$;
    \item extract a basis for $\mathcal{C}(A)$, $\mathcal{N}(A)$, and row space from REF;
    \item compute rank, nullity, and apply rank--nullity;
    \item express vectors uniquely in a basis (coordinates) and change basis;
    \item interpret dimension as degrees of freedom in computing systems.
\end{enumerate}

% ============================================================
\section{Linear Independence}
\label{sec:independence}
Picture three arrows in $\mathbb{R}^2$. The third must lie in the plane of the first two, so it is a combination of them --- redundant. Independence means no such redundancy.

\begin{definition}[Independence]
Vectors $\mathbf{v}_1,\dots,\mathbf{v}_k\in V$ are \emph{linearly independent} if
\[
  c_1\mathbf{v}_1+\cdots+c_k\mathbf{v}_k=\mathbf{0}\implies c_1=\cdots=c_k=0.
\]
Otherwise they are \emph{dependent}: some nontrivial combination gives $\mathbf{0}$ (equivalently, one vector is a combination of the others).
\end{definition}

Test: put vectors as columns of $A=[\mathbf{v}_1\cdots\mathbf{v}_k]$; independence $\iff\mathcal{N}(A)=\{\mathbf{0}\}\iff$ every column is a pivot column in REF.

\begin{figure}[h]
\centering
\begin{tikzpicture}[scale=0.82,every node/.style={font=\scriptsize}]
  % independent
  \begin{scope}
    \draw[->,thick] (-0.3,0)--(3.0,0); \draw[->,thick] (0,-0.3)--(0,2.6);
    \draw[->,very thick,blue!65!black] (0,0)--(2.2,0.5) node[midway,below] {$\mathbf{v}_1$};
    \draw[->,very thick,red!65!black] (0,0)--(0.6,2.0) node[midway,left] {$\mathbf{v}_2$};
    \fill[green!30,opacity=0.4] (0,0)--(2.2,0.5)--(2.8,2.5)--(0.6,2.0)--cycle;
    \node[font=\tiny,fill=white] at (1.4,1.2) {span fills plane};
    \node[font=\footnotesize\bfseries] at (1.4,-0.6) {Independent};
  \end{scope}
  % dependent
  \begin{scope}[xshift=4.5cm]
    \draw[->,thick] (-0.3,0)--(3.0,0); \draw[->,thick] (0,-0.3)--(0,2.6);
    \draw[->,very thick,blue!65!black] (0,0)--(2.0,0.5) node[midway,below] {$\mathbf{v}_1$};
    \draw[->,very thick,red!65!black] (0,0)--(1.0,0.25) node[midway,below] {$\mathbf{v}_2=0.5\mathbf{v}_1$};
    \draw[->,very thick,orange!75!black] (0,0)--(0.8,1.5) node[midway,left] {$\mathbf{v}_3$};
    \draw[dashed,gray] (0.8,1.5)--(1.6,0.4);
    \node[font=\tiny,fill=white] at (1.5,1.0) {$\mathbf{v}_3$ in span};
    \node[font=\footnotesize\bfseries] at (1.4,-0.6) {Dependent};
  \end{scope}
  \begin{scope}[xshift=9.0cm]
    \draw[->,thick] (-0.3,0)--(3.0,0); \draw[->,thick] (0,-0.3)--(0,2.6);
    \draw[->,very thick,blue!65!black] (0,0)--(2.0,0.6) node[midway,below] {$\mathbf{v}_1$};
    \draw[->,very thick,red!65!black] (0,0)--(0.7,1.9) node[midway,left] {$\mathbf{v}_2$};
    \draw[->,very thick,green!55!black] (0,0)--(2.7,2.5) node[midway,above right] {$\mathbf{v}_3=\mathbf{v}_1+\mathbf{v}_2$};
    \draw[dashed,gray] (2.0,0.6)--(2.7,2.5)--(0.7,1.9);
    \node[font=\footnotesize\bfseries] at (1.4,-0.6) {Dependent};
  \end{scope}
\end{tikzpicture}
\caption{Left: independent --- new direction each time. Middle/Right: dependent --- one vector lies in the span of the others.}
\label{fig:independence}
\end{figure}

\begin{example}[Polynomials]
In $\mathcal{P}_2$, $\{1,\;x,\;x^2\}$ is independent; $\{1,\;x,\;1+x\}$ is dependent since $(1+x)-1-x=0$. The test is the same: nontrivial coefficients making the zero polynomial.
\end{example}

\section{Span, Basis, Dimension}
\label{sec:basis}
\begin{definition}[Span and basis]
$\operatorname{span}\{\mathbf{v}_1,\dots,\mathbf{v}_k\}$ is the set of all combinations $\sum c_i\mathbf{v}_i$. A \emph{basis} of $V$ is a set that (i) is independent and (ii) spans $V$. The \emph{dimension} $\dim V$ is the number of vectors in any basis (proved well-defined below).
\end{definition}

Bases are not unique, but size is. Coordinates are unique once a basis is fixed.

\begin{theorem}[Dimension well-defined]
Any two bases of a finite-dimensional $V$ have the same cardinality. Any independent set can be extended to a basis; any spanning set contains a basis.
\end{theorem}
\begin{proof}[Sketch]
If $B$ is a basis ($n$ vectors) and $S$ is independent ($k$ vectors), each $\mathbf{s}\in S$ is a unique combination of $B$ --- the coordinate matrix is $n\times k$ with independent columns, so $k\le n$. Symmetrically $n\le k$ when both are bases.
\end{proof}

\begin{figure}[h]
\centering
\begin{tikzpicture}[scale=0.82,every node/.style={font=\scriptsize}]
  \draw[rounded corners=4pt,fill=blue!7,draw=blue!40!black,thick] (0,0) ellipse (2.6 and 1.5);
  \node[font=\footnotesize\bfseries] at (0,1.8) {Vector space $V$, $\dim V=n$};
  % independent set
  \node[draw,fill=green!15,rounded corners=3pt,inner sep=3pt,align=center] at (-1.4,0.4) {independent\\$k<n$: extend};
  \draw[->,thick,green!55!black] (-0.7,0.1)--(0.2,0.0) node[midway,above,font=\tiny] {add vectors};
  % spanning set
  \node[draw,fill=orange!15,rounded corners=3pt,inner sep=3pt,align=center] at (1.4,0.4) {spanning\\$k>n$: prune};
  \draw[->,thick,orange!70!black] (0.7,0.0)--(-0.15,-0.05) node[midway,below,font=\tiny] {remove redundancy};
  % basis in middle
  \node[draw,fill=violet!15,rounded corners=3pt,thick,inner sep=4pt,align=center] at (0,-0.9) {basis: $n$ vectors\\spans + independent};
  \draw[->,thick,violet!60!black] (0,-0.35)--(0,-0.55);
  \node[font=\tiny,fill=white] at (0,-1.5) {coordinates unique};
\end{tikzpicture}
\caption{A basis is a minimal spanning set and a maximal independent set; its size is the dimension.}
\label{fig:basis-dimension}
\end{figure}

\section{Bases for the Four Subspaces; Rank--Nullity}
\label{sec:rank-nullity}
For $A\in\mathbb{R}^{m\times n}$ with REF $R$:
\begin{itemize}
    \item $\mathcal{C}(A)$: basis = columns of $A$ at pivot positions of $R$ (not columns of $R$).
    \item $\mathcal{N}(A)$: free variables give one basis vector each (set one free $=1$, rest $0$, solve for pivots).
    \item Row space: nonzero rows of $R$ (or $R^T$ columns).
\end{itemize}
Let $r=\operatorname{rank}(A)=$ number of pivots $=\dim\mathcal{C}(A)=\dim\text{row space}$.

\begin{theorem}[Rank--nullity]
For $A\in\mathbb{R}^{m\times n}$,
\[
  \operatorname{rank}(A)+\operatorname{nullity}(A)=n,\qquad \operatorname{nullity}(A)=\dim\mathcal{N}(A).
\]
Also $\dim\mathcal{N}(A^T)=m-r$. So $r\le\min(m,n)$.
\end{theorem}

\begin{figure}[h]
\centering
\begin{tikzpicture}[scale=0.80,every node/.style={font=\scriptsize}]
  \node[draw,thick,fill=blue!10,minimum width=2.2cm,minimum height=1.4cm,rounded corners=3pt,align=center] (Rn) at (0,0) {$\mathbb{R}^n$\\[2pt] $\dim=n$};
  \node[draw,thick,fill=red!10,minimum width=2.2cm,minimum height=1.4cm,rounded corners=3pt,align=center] (Rm) at (6.0,0) {$\mathbb{R}^m$\\[2pt] $\dim=m$};
  \draw[->,very thick,violet!60!black] (1.3,0.3)--(4.7,0.3) node[midway,above] {$A$};
  \draw[->,very thick,orange!70!black] (4.7,-0.3)--(1.3,-0.3) node[midway,below] {$A^T$};
  \fill[green!30,opacity=0.5] (-0.6,-0.4) rectangle (0.6,0.2) node[midway,font=\tiny] {};
  \node[font=\tiny,fill=white] at (0,-0.75) {$\mathcal{N}(A)$, $\dim=n-r$};
  \fill[blue!30,opacity=0.5] (-0.9,0.25) rectangle (0.9,0.55);
  \node[font=\tiny] at (0,0.85) {row space, $\dim=r$};
  \fill[blue!30,opacity=0.5] (5.1,0.25) rectangle (6.9,0.55);
  \node[font=\tiny] at (6.0,0.85) {col space, $\dim=r$};
  \fill[red!30,opacity=0.5] (5.4,-0.4) rectangle (6.6,0.2);
  \node[font=\tiny,fill=white] at (6.0,-0.75) {$\mathcal{N}(A^T)$, $\dim=m-r$};
\end{tikzpicture}
\caption{The four subspaces and their dimensions. The row and column spaces both have dimension $r=\operatorname{rank}(A)$.}
\label{fig:four-dims}
\end{figure}

\begin{example}[Rank--nullity in action]
$A=\begin{bmatrix}1&2&1&0\\2&4&3&1\\1&2&2&1\end{bmatrix}\xrightarrow{\text{REF}}
\begin{bmatrix}1&2&1&0\\0&0&1&1\\0&0&0&0\end{bmatrix}$. Pivots in cols $1,3$ so $r=2$, nullity $=4-2=2$. $\mathcal{C}(A)$ basis: cols $1,3$ of $A$; $\mathcal{N}(A)$ basis: from $x_2,x_4$ free.
\end{example}

\begin{example}[Coordinates and change of basis]
Fix basis $B=\{\mathbf{b}_1,\mathbf{b}_2\}$. Then $[\mathbf{v}]_B$ solves $[\mathbf{b}_1\ \mathbf{b}_2][\mathbf{v}]_B=\mathbf{v}$. If $C$ is another basis, $[\mathbf{v}]_C=P_{C\gets B}[\mathbf{v}]_B$ where $P$ has columns $[\mathbf{b}_j]_C$.
\begin{lstlisting}[language=Python]
import numpy as np
B = np.array([[1.,1.],[0.,1.]]).T  # columns are b1,b2 -- careful with layout
# Actually B_cols = [[1,0],[1,1]] as columns
B_cols = np.array([[1.,1.],[0.,1.]])  # col j = B_cols[:,j]
v = np.array([3.,2.])
coords = np.linalg.solve(B_cols, v)
print(coords)  # [1., 2.] meaning v = 1*b1 + 2*b2
print(B_cols @ coords)  # [3., 2.] check
\end{lstlisting}
\end{example}

% ============================================================

% --- Additional depth: Exchange lemma, dimension theorem, and matrix rank ---
\section{Exchange and Dimension Uniqueness}
\label{sec:exchange}
\begin{lemma}[Steinitz exchange]
If $\{\mathbf{v}_1,\dots,\mathbf{v}_n\}$ spans $V$ and $\{\mathbf{u}_1,\dots,\mathbf{u}_k\}$ is independent, then $k\le n$ and the span can exchange $k$ of the $\mathbf{v}_i$ for the $\mathbf{u}_j$ and still span.
\end{lemma}
Corollary: all bases have the same size; any independent set extends to a basis (keep adding vectors outside the current span until spanning); any spanning set prunes to a basis (drop dependent vectors one by one). This justifies calling $\dim V$ well-defined.

\begin{example}[Extending in $\mathcal{P}_2$]
$\{\mathbf{v}_1=1+x\}$ independent in $\mathcal{P}_2$ ($\dim3$) extends to a basis, e.g.\ $\{1+x,\;x,\;x^2\}$ or $\{1+x,\;1-x,\;x^2\}$. Infinitely many bases; dimension is the invariant.
\end{example}

\section{Matrix Rank = Row Rank = Column Rank}
\label{sec:rank-equal}
A nontrivial theorem: $\dim\mathcal{C}(A)=\dim\mathcal{C}(A^T)$ --- column rank equals row rank, both $=r$. So REF pivot count is intrinsic, not an artefact. Proof sketch: show $\mathcal{N}(A)$ is orthogonal complement of row space (Ch.~6); then rank--nullity gives equality.

\section{Change of Basis and Similarity Preview}
\label{sec:change-basis-detail}
If $B$ and $C$ are bases, $P_{C\gets B}$ (columns $[\mathbf{b}_j]_C$) converts $[\mathbf{v}]_B$ to $[\mathbf{v}]_C$. For a linear operator $T:V\to V$, its matrices satisfy $[T]_C=P[T]_BP^{-1}$ --- similarity. Determinant, trace, eigenvalues are similarity invariants (Ch.~7--8). This is why diagonalization $A=PDP^{-1}$ is a change to the eigenbasis where $T$ acts diagonally.


\section{Chapter Notes}
Grassmann (1844) first framed dimension this way. Rank--nullity is the ``conservation of dimensions'' for $A:\mathbb{R}^n\to\mathbb{R}^m$. Computing rank via SVD (Ch.~6/8 notes) is numerically stable; REF rank is exact-arithmetic.

% ============================================================
\section{Exercises}
\label{sec:ch05-exercises}
\begin{enumerate}[label=E5.\arabic*, leftmargin=*, itemsep=0.8em]
\item \textbf{Test.} Are $\mathbf{v}_1=[1,0,1]^T$, $\mathbf{v}_2=[0,1,1]^T$, $\mathbf{v}_3=[1,1,2]^T$ independent? If dependent, write one as a combination of the others.
\item \textbf{Bases from REF.} For $A=\begin{bmatrix}1&2&0&1\\0&0&1&2\\1&2&1&3\end{bmatrix}$, find bases for $\mathcal{C}(A)$, $\mathcal{N}(A)$, and row space. Verify rank$+$nullity$=n$ and $r\le\min(m,n)$.
\item \textbf{Coordinates.} Let $B=\{[1,1]^T,[1,-1]^T\}$ be a basis of $\mathbb{R}^2$. Find $[\mathbf{v}]_B$ for $\mathbf{v}=[3,1]^T$. If $C$ is the standard basis, write $P_{C\gets B}$ and $P_{B\gets C}$.
\item \textbf{Dimension bound.} Prove any $k>n$ vectors in $\mathbb{R}^n$ are dependent. Deduce $\dim\mathbb{R}^n=n$ via the standard basis.
\item \textbf{Coding.} Write \texttt{col\_basis(A)} returning indices of pivot columns (tolerance $10^{-10}$) and a basis for $\mathcal{N}(A)$ via SVD or REF. Compare with \texttt{scipy.linalg.null\_space} on random $5\times7$ matrices.
\end{enumerate}
% End of Chapter 5

% SC1004 Chapter 6: Orthogonality & Least Squares (0->100 ground-up)
\chapter{Orthogonality and Least Squares}
\label{ch:orthogonality}

\begin{quote}
\emph{``Orthogonality is geometry's way of saying `no shadow'.''} When a system $A\mathbf{x}=\mathbf{b}$ has no solution, the closest we can do is the shadow of $\mathbf{b}$ on the column space --- that shadow solves the least-squares problem.
\end{quote}

\noindent\fbox{\parbox{\dimexpr\linewidth-2\fboxsep-2\fboxrule}{%
\textbf{From zero: right angles before residuals.} We generalise $\mathbf{u}\cdot\mathbf{v}=0$ from $\mathbb{R}^2$ to $\mathbb{R}^n$, build orthogonal projections, then solve $A\mathbf{x}\approx\mathbf{b}$ when equality is impossible --- the workhorse of regression and data fitting.}}
\vspace{0.4em}

\section*{Learning Outcomes}
After this chapter you will be able to:
\begin{enumerate}[label=(L\arabic*)]
    \item test orthogonality and orthonormality of sets;
    \item compute orthogonal projections onto subspaces;
    \item run Gram--Schmidt to build orthonormal bases and $QR$;
    \item formulate and solve least-squares via normal equations and $QR$;
    \item apply least squares to line fitting and explain residuals.
\end{enumerate}

% ============================================================
\section{Orthogonality in \texorpdfstring{$\mathbb{R}^n$}{Rn}}
\label{sec:ortho-def}
Recall $\mathbf{u}\perp\mathbf{v}\iff\mathbf{u}\cdot\mathbf{v}=0$ (Ch.~1). A set $\{\mathbf{q}_1,\dots,\mathbf{q}_k\}$ is \emph{orthogonal} if each pair is orthogonal; \emph{orthonormal} if in addition $\|\mathbf{q}_i\|=1$. Orthonormal sets are automatically independent --- no one can be built from perpendicular directions.

\begin{figure}[h]
\centering
\begin{tikzpicture}[scale=0.85,every node/.style={font=\scriptsize}]
  \draw[->,thick] (-0.3,0)--(3.2,0) node[right] {$x$};
  \draw[->,thick] (0,-0.3)--(0,3.0) node[above] {$y$};
  \draw[->,very thick,blue!65!black] (0,0)--(2.4,0.6) node[midway,below] {$\mathbf{q}_1$};
  \draw[->,very thick,red!65!black] (0,0)--(-0.5,2.0) node[midway,left] {$\mathbf{q}_2$};
  \draw[thick,green!60!black] (0.35,0.09)--(0.26,0.44)--(-0.09,0.35);
  \node[font=\tiny,fill=white] at (0.45,0.4) {$90^\circ$};
  \draw[dashed,gray] (2.4,0.6)--(1.9,2.6)--(-0.5,2.0);
  \fill[blue!8,opacity=0.4] (0,0)--(2.4,0.6)--(1.9,2.6)--(-0.5,2.0)--cycle;
  \node[font=\tiny,align=center] at (1.0,1.4) {rectangle if\\orthonormal + equal length};
  \node[font=\footnotesize\bfseries] at (1.4,-0.6) {$\mathbf{q}_1\perp\mathbf{q}_2$ in $\mathbb{R}^2$};

  \begin{scope}[xshift=5.0cm]
    \draw[->,thick] (-0.3,0)--(3.2,0); \draw[->,thick] (0,-0.3)--(0,3.0);
    \draw[->,very thick,blue!65!black] (0,0)--(1.8,0.5) node[midway,below] {$\mathbf{a}_1$};
    \draw[->,very thick,red!65!black] (0,0)--(0.9,1.8) node[midway,left] {$\mathbf{a}_2$};
    \draw[thick,orange!70!black] (0.4,0.11)--(0.33,0.36)--(0.08,0.29);
    \node[font=\tiny,fill=white] at (0.7,0.6) {not $90^\circ$};
    \draw[->,thick,violet!70!black,dashed] (0.9,1.8)--(1.45,0.4) node[midway,right,font=\tiny] {overlap};
    \node[font=\footnotesize\bfseries] at (1.4,-0.6) {Non-orthogonal};
  \end{scope}
\end{tikzpicture}
\caption{Left: orthogonal vectors meet at $90^\circ$; orthonormal adds unit length. Right: general non-orthogonal pair (Gram--Schmidt will fix this).}
\label{fig:orthogonality}
\end{figure}

If $Q=[\mathbf{q}_1\cdots\mathbf{q}_k]$ has orthonormal columns, then $Q^TQ=I_k$ --- columns are perfectly conditioned; $Q$ preserves lengths: $\|Q\mathbf{x}\|=\|\mathbf{x}\|$.

\section{Projections onto Subspaces}
\label{sec:projections}
Let $W=\operatorname{span}\{\mathbf{a}_1,\dots,\mathbf{a}_k\}\subseteq\mathbb{R}^n$. For $\mathbf{b}\in\mathbb{R}^n$, the \emph{projection} $\hat{\mathbf{b}}=\operatorname{proj}_W\mathbf{b}$ is the unique $\mathbf{w}\in W$ closest to $\mathbf{b}$; the residual $\mathbf{b}-\hat{\mathbf{b}}$ is orthogonal to $W$.

When $W$ has orthonormal basis $Q$,
\[
  \operatorname{proj}_W\mathbf{b}=QQ^T\mathbf{b}=\sum_{i=1}^{k}(\mathbf{q}_i^T\mathbf{b})\,\mathbf{q}_i.
\]
In general with $A=[\mathbf{a}_1\cdots\mathbf{a}_k]$ (full column rank),
\[
  \operatorname{proj}_W\mathbf{b}=A(A^TA)^{-1}A^T\mathbf{b}.
\]
The matrix $P=A(A^TA)^{-1}A^T$ satisfies $P^2=P$, $P^T=P$ --- the orthogonal projector onto $\mathcal{C}(A)$.

\begin{figure}[h]
\centering
\begin{tikzpicture}[scale=0.82,every node/.style={font=\scriptsize}]
  \fill[blue!10] (-1.2,-0.3)--(3.4,0.2)--(2.2,2.8)--(-1.4,2.3)--cycle;
  \draw[very thick,blue!50!black] (-1.2,-0.3)--(3.4,0.2)--(2.2,2.8)--(-1.4,2.3)--cycle;
  \node[font=\footnotesize\bfseries,blue!50!black] at (3.0,2.6) {$W=\mathcal{C}(A)$};
  \fill[red!70!black] (1.0,3.2) circle (2.5pt) node[above] {$\mathbf{b}$};
  \fill[green!60!black] (0.85,1.0) circle (2.5pt) node[below right] {$\hat{\mathbf{b}}=\operatorname{proj}_W\mathbf{b}$};
  \draw[->,very thick,red!60!black,dashed] (0.85,1.0)--(1.0,3.2) node[midway,right] {$\mathbf{b}-\hat{\mathbf{b}}\perp W$};
  \draw[thick,green!60!black] (0.70,1.18)--(0.85,1.32)--(0.99,1.15);
  \draw[->,thick,violet!70!black] (1.0,3.2)--(0.85,1.0) node[midway,left,font=\tiny] {shortest};
  \node[font=\tiny,fill=white] at (0.9,-0.5) {Residual orthogonal to $W$; $\hat{\mathbf{b}}$ is the closest point in $W$ to $\mathbf{b}$};
\end{tikzpicture}
\caption{Orthogonal projection: $\hat{\mathbf{b}}$ is the closest point in $W$ to $\mathbf{b}$; the error $\mathbf{b}-\hat{\mathbf{b}}$ is perpendicular to $W$.}
\label{fig:projection-subspace}
\end{figure}

\section{Gram--Schmidt and QR}
\label{sec:gram-schmidt}
Gram--Schmidt turns any independent $\mathbf{a}_1,\dots,\mathbf{a}_k$ into orthonormal $\mathbf{q}_1,\dots,\mathbf{q}_k$ spanning the same nested subspaces:

\begin{align}
  \mathbf{u}_1&=\mathbf{a}_1,\quad \mathbf{q}_1=\mathbf{u}_1/\|\mathbf{u}_1\|,\\
  \mathbf{u}_2&=\mathbf{a}_2-(\mathbf{q}_1^T\mathbf{a}_2)\mathbf{q}_1,\quad \mathbf{q}_2=\mathbf{u}_2/\|\mathbf{u}_2\|,\\
  &\ \vdots\\
  \mathbf{u}_k&=\mathbf{a}_k-\sum_{i<k}(\mathbf{q}_i^T\mathbf{a}_k)\mathbf{q}_i.
\end{align}
Each step subtracts the shadows already accounted for. With $A=QR$ ($Q$ orthonormal columns, $R$ upper-triangular), $R=Q^TA$ records the coefficients.

\begin{figure}[h]
\centering
\begin{tikzpicture}[scale=0.82,every node/.style={font=\scriptsize}]
  \draw[->,thick] (-0.3,0)--(3.5,0); \draw[->,thick] (0,-0.3)--(0,3.0);
  \draw[->,very thick,blue!65!black] (0,0)--(2.2,0.6) node[midway,below] {$\mathbf{a}_1\to\mathbf{q}_1$};
  \draw[->,very thick,red!65!black] (0,0)--(1.0,2.2) node[midway,left] {$\mathbf{a}_2$};
  \draw[->,thick,orange!70!black,dashed] (1.0,2.2)--(1.72,0.48) node[midway,right,font=\tiny] {subtract proj};
  \draw[->,very thick,green!60!black] (0,0)--(0.62,1.58) node[midway,left] {$\mathbf{u}_2\to\mathbf{q}_2$};
  \draw[thick,green!60!black] (1.50,0.41)--(1.57,0.66)--(1.32,0.73);
  \node[font=\tiny,fill=white] at (1.2,1.0) {$\mathbf{u}_2\perp\mathbf{q}_1$};
  \node[font=\footnotesize\bfseries] at (1.6,-0.6) {Gram--Schmidt step 2};
\end{tikzpicture}
\caption{Gram--Schmidt: subtract from $\mathbf{a}_2$ its projection on $\mathbf{q}_1$; what remains is orthogonal to $\mathbf{q}_1$. Normalise to get $\mathbf{q}_2$.}
\label{fig:gram-schmidt}
\end{figure}

\begin{example}[QR by hand ($2\times2$)]
$A=\begin{bmatrix}1&1\\1&0\\0&1\end{bmatrix}$ ($3\times2$): $\mathbf{q}_1=[1,1,0]^T/\sqrt2$, $\mathbf{u}_2=[1,0,1]^T-(\mathbf{q}_1^T[1,0,1]^T)\mathbf{q}_1=[1/2,-1/2,1]^T$, $\mathbf{q}_2=\mathbf{u}_2/\|\mathbf{u}_2\|$. Then $R=Q^TA$ is $2\times2$ upper-triangular.
\end{example}

\section{Least Squares}
\label{sec:least-squares}
When $A\mathbf{x}=\mathbf{b}$ has no solution ($\mathbf{b}\notin\mathcal{C}(A)$), minimise the squared error:
\[
  \min_{\mathbf{x}}\|A\mathbf{x}-\mathbf{b}\|^2.
\]
The minimiser $\hat{\mathbf{x}}$ makes $A\hat{\mathbf{x}}=\operatorname{proj}_{\mathcal{C}(A)}\mathbf{b}$, so the residual is orthogonal to $\mathcal{C}(A)$: $A^T(\mathbf{b}-A\hat{\mathbf{x}})=\mathbf{0}$.

\begin{theorem}[Normal equations]
$\hat{\mathbf{x}}$ is a least-squares solution iff
\[
  A^TA\hat{\mathbf{x}}=A^T\mathbf{b}.
\]
If $A$ has independent columns, $A^TA$ is invertible and $\hat{\mathbf{x}}=(A^TA)^{-1}A^T\mathbf{b}$ is unique. With $A=QR$, $\hat{\mathbf{x}}=R^{-1}Q^T\mathbf{b}$ (more stable).
\end{theorem}

\begin{example}[Line fitting]
Fit $y=c+dt$ through $(t_i,y_i)$: $A=\begin{bmatrix}1&t_1\\\vdots&\vdots\\1&t_m\end{bmatrix}$, $\mathbf{x}=[c,d]^T$. Then $A^TA=\begin{bmatrix}m&\sum t_i\\\sum t_i&\sum t_i^2\end{bmatrix}$, $A^T\mathbf{b}=[\sum y_i,\sum t_iy_i]^T$. Solving gives the regression line; $\|\mathbf{b}-A\hat{\mathbf{x}}\|$ is the residual.
\begin{lstlisting}[language=Python]
import numpy as np
t = np.array([0.,1.,2.,3.])
y = np.array([1.1,1.9,3.2,3.8])
A = np.column_stack([np.ones_like(t), t])
x_hat, residuals, rank, s = np.linalg.lstsq(A, y, rcond=None)
print(x_hat)  # [c, d] intercept + slope
# Compare with normal equations:
print(np.linalg.solve(A.T@A, A.T@y))
\end{lstlisting}
\end{example}

% ============================================================

% --- Additional depth: Orthogonal complements and best approximation ---
\section{Orthogonal Complements and the Four Subspaces}
\label{sec:ortho-complement}
For $W\subseteq\mathbb{R}^n$, $W^\perp=\{\mathbf{x}:\mathbf{x}\perp\mathbf{w}\;\forall\mathbf{w}\in W\}$ is a subspace with $\dim W+\dim W^\perp=n$ and $(W^\perp)^\perp=W$. Fundamental theorem:
\[
  \mathcal{C}(A)^\perp=\mathcal{N}(A^T),\qquad \mathcal{C}(A^T)^\perp=\mathcal{N}(A).
\]
So $\mathbb{R}^n=\mathcal{C}(A^T)\oplus\mathcal{N}(A)$ and $\mathbb{R}^m=\mathcal{C}(A)\oplus\mathcal{N}(A^T)$ as orthogonal direct sums. The least-squares residual lives in $\mathcal{N}(A^T)$.

\section{Best Approximation and the Projection Theorem}
\label{sec:best-approx}
\begin{theorem}[Projection theorem]
Let $W\subseteq\mathbb{R}^n$ be a subspace. For any $\mathbf{b}$, there is a unique $\hat{\mathbf{b}}\in W$ minimising $\|\mathbf{b}-\hat{\mathbf{b}}\|$, namely $\hat{\mathbf{b}}=\operatorname{proj}_W\mathbf{b}$. Moreover $\mathbf{b}-\hat{\mathbf{b}}\in W^\perp$.
\end{theorem}
This underpins all of least squares: $A\hat{\mathbf{x}}$ is the best approximation to $\mathbf{b}$ from $\mathcal{C}(A)$. Error $\|\mathbf{b}-A\hat{\mathbf{x}}\|$ is the distance from $\mathbf{b}$ to the subspace.

\section{Numerical Stability: QR vs.\ Normal Equations}
\label{sec:stability}
Forming $A^TA$ squares the condition number $\kappa(A^TA)=\kappa(A)^2$, so normal equations lose half the digits when $A$ is ill-conditioned. $QR$ avoids this: $Q$ is perfectly conditioned ($\kappa(Q)=1$), so solving $R\hat{\mathbf{x}}=Q^T\mathbf{b}$ is stable. Modified Gram--Schmidt and Householder $QR$ are preferred; NumPy's \texttt{lstsq} uses SVD/Householder internally.



\begin{example}[Polynomial least squares]
Fit a parabola $y=c_0+c_1t+c_2t^2$ to $m$ points: $A$ is $m\times3$ Vandermonde $[1\ t_i\ t_i^2]$. Columns become nearly dependent for large $t_i$ (ill-conditioned); centering $t_i$ or using orthogonal polynomials (Legendre) fixes this. \texttt{np.polyfit} does a scaled $QR$ internally for this reason.
\end{example}

\begin{remark}[When $A^TA$ is singular]
If $A$ has dependent columns, $A^TA$ singular and normal equations have infinitely many solutions. The minimal-norm least-squares solution is $A^\dagger\mathbf{b}$ where $A^\dagger=V\Sigma^\dagger U^T$ is the pseudoinverse (SVD, Ch.~8). NumPy's \texttt{lstsq} returns this by default.
\end{remark}



\begin{remark}[Geometric meaning of $Q^TQ=I$]
Columns $\mathbf{q}_i$ orthonormal $\iff$ $Q$ is an isometry on its column space: angles and lengths preserved. In graphics, $Q$ is a rotation/reflection; in data science, $Q$ from $QR$ is the ``well-conditioned'' part of $A$. Never form $Q$ explicitly when $m\gg n$ --- Householder applies $Q^T\mathbf{b}$ implicitly in $O(mn)$.
\end{remark}


\section{Chapter Notes}
Least squares (Legendre, Gauss, 1805--1809) predates matrices; Gauss used it to predict Ceres' orbit. $QR$ via Gram--Schmidt or Householder reflections is the stable solver (LAPACK's \texttt{gels}); forming $A^TA$ squares the condition number and is avoided for ill-conditioned $A$.

% ============================================================
\section{Exercises}
\label{sec:ch06-exercises}
\begin{enumerate}[label=E6.\arabic*, leftmargin=*, itemsep=0.8em]
\item \textbf{Orthonormal check.} Is $Q=\frac12\begin{bmatrix}1&1&1&1\\1&-1&1&-1\\1&1&-1&-1\\1&-1&-1&1\end{bmatrix}$ orthogonal ($Q^TQ=I$)? What is $Q\mathbf{x}$ geometrically?
\item \textbf{Gram--Schmidt.} Apply Gram--Schmidt to $\mathbf{a}_1=[1,1,0]^T$, $\mathbf{a}_2=[1,0,1]^T$, $\mathbf{a}_3=[0,1,1]^T$ in $\mathbb{R}^3$. Write $A=QR$ and verify $Q^TQ=I$, $A=QR$.
\item \textbf{Projection matrix.} For $A=[1,1,1]^T$ (span of all-ones vector in $\mathbb{R}^3$), compute $P=A(A^TA)^{-1}A^T$. Show $P^2=P$, $P^T=P$, and $P\mathbf{b}$ is the mean of entries times $\mathbf{1}$.
\item \textbf{Least squares.} Fit $y=c+dt$ to $(0,1),(1,2),(2,2),(3,4)$ via normal equations. Compute $\hat{\mathbf{x}}$ and $\|\mathbf{b}-A\hat{\mathbf{x}}\|$. Verify $A^T(\mathbf{b}-A\hat{\mathbf{x}})=\mathbf{0}$.
\item \textbf{Coding.} Generate $m=100$ noisy points $y=2+3t+\mathcal{N}(0,0.5)$, $t\in[0,5]$. Solve least squares via (a) normal equations and (b) \texttt{np.linalg.lstsq} ($QR$). Compare $\hat{\mathbf{x}}$ and residuals; which is more stable if you add an outlier?
\end{enumerate}
% End of Chapter 6

% SC1004 Chapter 7: Determinants (0->100 ground-up)
\chapter{Determinants}
\label{ch:determinants}

\begin{quote}
\emph{``The determinant is the signed volume of the box the matrix builds.''} Zero means the box is flat --- the matrix squashes space and cannot be undone.
\end{quote}

\noindent\fbox{\parbox{\dimexpr\linewidth-2\fboxsep-2\fboxrule}{%
\textbf{From zero: area before algebra.} We start with area of a parallelogram in $\mathbb{R}^2$ ($ad-bc$), grow to volume in $\mathbb{R}^3$, then axiomatise to $n\times n$ via row operations. Every property is first a picture.}}
\vspace{0.4em}

\section*{Learning Outcomes}
After this chapter you will be able to:
\begin{enumerate}[label=(L\arabic*)]
    \item compute $2\times2$ and $3\times3$ determinants and interpret them as area/volume;
    \item apply row-operation rules and multiplicativity $\det(AB)=\det A\det B$;
    \item use $\det A\neq0$ to test invertibility and find $\det(A^{-1})$;
    \item expand by cofactors and relate to the adjugate;
    \item connect $\det$ to eigenvalues and to Cramer's rule.
\end{enumerate}

% ============================================================
\section{Area, Volume, and the \texorpdfstring{$2\times2$}{2x2} Determinant}
\label{sec:det2x2}
Take columns $\mathbf{a}_1,\mathbf{a}_2\in\mathbb{R}^2$ as edges from the origin. They span a parallelogram. Its signed area is
\[
  \det\begin{bmatrix}a&b\\c&d\end{bmatrix}=ad-bc.
\]
Sign encodes orientation: $\det>0$ means $\mathbf{a}_2$ is counter-clockwise from $\mathbf{a}_1$.

\begin{figure}[h]
\centering
\begin{tikzpicture}[scale=0.85,every node/.style={font=\scriptsize}]
  \draw[->,thick] (-0.4,0)--(3.5,0) node[right] {$x$};
  \draw[->,thick] (0,-0.4)--(0,3.0) node[above] {$y$};
  \fill[blue!15] (0,0)--(2.2,0.5)--(2.9,2.3)--(0.7,1.8)--cycle;
  \draw[very thick,blue!60!black] (0,0)--(2.2,0.5) node[midway,below] {$\mathbf{a}_1$};
  \draw[very thick,red!60!black] (0,0)--(0.7,1.8) node[midway,left] {$\mathbf{a}_2$};
  \draw[thick,orange!60!black,dashed] (2.2,0.5)--(2.9,2.3)--(0.7,1.8);
  \node[font=\tiny,fill=white] at (1.45,1.15) {area $=|\det A|$};
  \draw[->,thick,green!55!black] (1.1,0.25) arc (12:68:1.1) node[midway,right,font=\tiny] {$\det>0$};
  \node[font=\footnotesize\bfseries] at (1.45,-0.6) {Parallelogram area $=|ad-bc|$};

  \begin{scope}[xshift=5.5cm]
    \draw[->,thick] (-0.4,0)--(3.2,0); \draw[->,thick] (0,-0.4)--(0,3.0);
    \fill[red!15] (0,0)--(2.0,0.4)--(2.0,0.4)--cycle;
    \draw[very thick,blue!60!black] (0,0)--(2.0,0.4) node[midway,below] {$\mathbf{a}_1$};
    \draw[very thick,red!60!black] (0,0)--(1.0,0.2) node[midway,above] {$\mathbf{a}_2=0.5\mathbf{a}_1$};
    \node[font=\tiny,fill=white] at (1.0,1.0) {area $0$};
    \node[font=\tiny,fill=white] at (1.0,1.35) {collinear $\Rightarrow\det=0$};
    \node[font=\footnotesize\bfseries] at (1.4,-0.6) {Degenerate: $\det=0$};
  \end{scope}
\end{tikzpicture}
\caption{Left: $\det A$ is the signed area of the column parallelogram. Right: dependent columns give zero area --- the box is flat.}
\label{fig:det-area}
\end{figure}

In $\mathbb{R}^3$, $\det[\mathbf{a}_1\ \mathbf{a}_2\ \mathbf{a}_3]$ is the signed volume of the parallelepiped. The pattern: determinant measures how the matrix scales volumes.

\section{Axioms and Row-Operation Rules}
\label{sec:det-axioms}
For $n\times n$, $\det$ is the unique function $\mathbb{R}^{n\times n}\to\mathbb{R}$ that is (i) multilinear in rows, (ii) alternating (zero if two rows equal), (iii) $\det I_n=1$. From this:

\begin{figure}[h]
\centering
\begin{tikzpicture}[scale=0.78,every node/.style={font=\scriptsize}]
  \node[draw,thick,rounded corners=3pt,fill=blue!10,minimum width=3.8cm,minimum height=0.7cm] (r1) at (0,1.8) {Swap rows: $\det$ flips sign};
  \node[draw,thick,rounded corners=3pt,fill=green!10,minimum width=3.8cm,minimum height=0.7cm] (r2) at (0,0.9) {Scale row by $c$: $\det$ scales by $c$};
  \node[draw,thick,rounded corners=3pt,fill=orange!12,minimum width=3.8cm,minimum height=0.7cm] (r3) at (0,0.0) {$R_i\gets R_i+cR_j$: $\det$ unchanged};
  \node[draw,thick,rounded corners=3pt,fill=violet!10,minimum width=4.2cm,minimum height=0.7cm] (r4) at (0,-0.9) {Triangular $\det=$ product of diagonal};
  \draw[->,thick,blue!60!black] (r1)--(r2);
  \draw[->,thick,green!55!black] (r2)--(r3);
  \draw[->,thick,orange!70!black] (r3)--(r4);
  \node[font=\tiny,align=center] at (5.0,0.45) {So elimination\\computes $\det$ in\\$O(n^3)$: reduce\\to triangular,\\multiply diagonal\\(track swaps/scales).};
  \draw[->,thick,gray] (2.0,0.45)--(3.4,0.45);
\end{tikzpicture}
\caption{How row operations affect the determinant. Elimination to triangular form is the practical algorithm.}
\label{fig:det-rowops}
\end{figure}

Consequences:
\begin{itemize}
    \item $\det(AB)=\det A\cdot\det B$, $\det(A^T)=\det A$.
    \item $A$ invertible $\iff\det A\neq0$, and $\det(A^{-1})=1/\det A$.
    \item $\det(cA)=c^n\det A$ for $n\times n$ (each of $n$ rows scaled).
    \item Row operations as above; column operations obey the same rules ($\det A^T=\det A$).
\end{itemize}

\begin{example}[Via elimination]
$A=\begin{bmatrix}2&1&0\\0&3&1\\1&0&2\end{bmatrix}$. Swap $R_1\leftrightarrow R_3$ flips sign: $\det A=-\det\begin{bmatrix}1&0&2\\0&3&1\\2&1&0\end{bmatrix}$. Then $R_3\gets R_3-2R_1$: $-\det\begin{bmatrix}1&0&2\\0&3&1\\0&1&-4\end{bmatrix}=-1\cdot\det\begin{bmatrix}3&1\\1&-4\end{bmatrix}=-1\cdot(-13)=13$.
\end{example}

\section{Cofactor Expansion and Adjugate}
\label{sec:cofactor}
Delete row $i$, column $j$ from $A$ to get $M_{ij}$ ($n-1\times n-1$). The \emph{cofactor} is $C_{ij}=(-1)^{i+j}\det M_{ij}$. Then for any fixed $i$,
\[
  \det A=\sum_{j=1}^{n} a_{ij}C_{ij}\quad\text{(row $i$ expansion)},
\]
and similarly for any column. The \emph{adjugate} $\operatorname{adj}(A)=[C_{ji}]$ satisfies
\[
  A\cdot\operatorname{adj}(A)=\det(A)\,I_n,
\]
so if $\det A\neq0$, $A^{-1}=\operatorname{adj}(A)/\det A$. For $2\times2$, this is $\frac1{ad-bc}\begin{bmatrix}d&-b\\-c&a\end{bmatrix}$.

Cofactor expansion is $O(n!)$ --- never used for large $n$; elimination is $O(n^3)$.

\begin{figure}[h]
\centering
\begin{tikzpicture}[scale=0.80,every node/.style={font=\scriptsize}]
  \foreach \r in {0,...,2} \foreach \c in {0,...,2} {
    \pgfmathparse{(\r==1 && \c==1) ? 1 : 0}
    \ifnum\pgfmathresult=1
      \node[draw,thick,minimum width=0.85cm,minimum height=0.60cm,fill=red!25] at (\c*0.95,\r*-0.70) {$a_{22}$};
    \else
      \pgfmathparse{(\r==1 || \c==1) ? 1 : 0}
      \ifnum\pgfmathresult=1
        \node[draw,minimum width=0.85cm,minimum height=0.60cm,fill=gray!25] at (\c*0.95,\r*-0.70) {$\times$};
      \else
        \node[draw,minimum width=0.85cm,minimum height=0.60cm,fill=blue!10] at (\c*0.95,\r*-0.70) {$\cdot$};
      \fi
    \fi
  }
  \draw[very thick,red!60!black,rounded corners=2pt] (-0.52,0.35)--(2.35,0.35)--(2.35,-1.80)--(-0.52,-1.80)--cycle;
  \node[font=\tiny,fill=white] at (1.0,-2.15) {Delete row 2, col 2 $\to M_{22}$ (blue $2\times2$)};
  \node[font=\tiny,fill=white] at (4.2, -0.65) {$C_{22}=(-1)^{4}\det M_{22}$};
  \draw[->,thick,blue!60!black] (2.85,-0.65)--(3.45,-0.65);
\end{tikzpicture}
\caption{Cofactor $C_{ij}$: delete row $i$ and column $j$, take the determinant of the remaining $(n-1)\times(n-1)$ minor, multiply by $(-1)^{i+j}$.}
\label{fig:cofactor}
\end{figure}

\section{Determinant and Eigenvalues}
\label{sec:det-eigen-preview}
If $A\mathbf{v}=\lambda\mathbf{v}$, then $\det(A-\lambda I)=0$ (Ch.~8). Hence $\det A=\prod_i\lambda_i$ (product of eigenvalues) and $\operatorname{tr}A=\sum_i\lambda_i$. Cramer's rule $x_j=\det A_j/\det A$ (replace column $j$ of $A$ by $\mathbf{b}$) follows from the adjugate but is $O(n!)$ --- a theoretical tool, not an algorithm.

\begin{example}[Python --- determinants]
\begin{lstlisting}[language=Python]
import numpy as np
A = np.array([[2.,1.,0.],[0.,3.,1.],[1.,0.,2.]])
print(np.linalg.det(A))          # ~13.0
print(np.linalg.det(A.T))        # same
print(np.linalg.det(A @ A))      # det(A)^2
# Volume scaling: unit cube -> parallelepiped of volume |det A|
print(abs(np.linalg.det(A)))
\end{lstlisting}
\end{example}

% ============================================================

% --- Additional depth: Properties, geometry in n-dim, and volume scaling ---
\section{Multiplicativity and Invertibility Revisited}
\label{sec:mult-inv}
$\det(AB)=\det A\det B$ implies $\det(A^{-1})=1/\det A$ and $\det(PAP^{-1})=\det A$ (similarity preserves $\det$). Hence $\det$ is a basis-free invariant of the linear map $T:\mathbf{x}\mapsto A\mathbf{x}$: it is the volume scaling factor of $T$, independent of coordinates.

\begin{theorem}[Volume scaling]
For any measurable $S\subset\mathbb{R}^n$, $\operatorname{vol}(A(S))=|\det A|\cdot\operatorname{vol}(S)$ when $A$ is invertible. $|\det A|$ is the factor by which $A$ stretches $n$-dimensional volume; $\det A=0$ collapses volume to zero.
\end{theorem}
This is why $|\det J|$ appears as the Jacobian in change of variables for integrals.

\section{Determinant, Rank, and Eigenvalues}
\label{sec:det-rank-eigen}
$\det A\neq0\iff\operatorname{rank}A=n\iff0$ is not an eigenvalue. More generally, $\det A=\prod_{i=1}^n\lambda_i$ where $\lambda_i$ are eigenvalues (with algebraic multiplicity, possibly complex). So $|\det A|$ is the product of singular values too (Ch.~8 notes). For triangular $A$, $\det$ is the diagonal product --- eigenvalues sit on the diagonal.

\begin{example}[When $\det$ misleads]
$A=0.1\,I_{100}$ has $\det A=10^{-100}$ (tiny) but $A$ is perfectly invertible ($\kappa=1$). Determinant magnitude does not measure near-singularity; condition number does. Never test ``$\det\approx0$'' for invertibility in floating point.
\end{example}

\section{Cramer's Rule and Its Cost}
\label{sec:cramer-detail}
If $A$ invertible, $A^{-1}=\operatorname{adj}(A)/\det A$ and $x_j=\det A_j/\det A$ where $A_j$ replaces column $j$ by $\mathbf{b}$. Cost $O(n!)$ via cofactors vs.\ $O(n^3)$ via $LU$: for $n=20$, $20!\approx2.4\times10^{18}$ --- intractable. Cramer is a theoretical gem that proves $A^{-1}$ entries are rational functions of $A$'s entries.



\begin{remark}[Orientation and handedness]
In graphics, $\det>0$ means a right-handed coordinate frame stays right-handed after $A$; $\det<0$ means it flips (mirror). This is why $\det$ of a rotation is $+1$ and of a reflection is $-1$ --- the sign is the parity of the transformation.
\end{remark}

\begin{example}[Determinant and invertibility in code]
Never branch on \texttt{det(A)==0} in floating point; use \texttt{np.linalg.matrix\_rank} or \texttt{cond}. For $n=100$, round-off makes $\det$ under/overflow ($\sim10^{\pm100}$) while rank via SVD is reliable. Check \texttt{np.linalg.cond(A)} --- large means near-singular regardless of $\det$ scale.
\end{example}



\begin{example}[Triangular shortcut]
If $U$ is upper-triangular, $\det U=\prod u_{ii}$ immediately. So in $PA=LU$, $\det A=\pm\prod u_{ii}$ (sign $=$ parity of row swaps in $P$). This is the $O(n^3)$ algorithm; no cofactor needed. For $2\times2$, both routes agree: $\det\begin{bmatrix}a&b\\c&d\end{bmatrix}=ad-bc$.
\end{example}


\section{Chapter Notes}
Cauchy and Jacobi formalised $\det(AB)=\det A\det B$; Laplace gave the expansion. Modern libraries compute $\det$ via $LU$ with partial pivoting and track sign; for large $n$, $\log|\det|$ is returned to avoid overflow. The volume picture generalises to $|\det|$ as the Jacobian in multivariate calculus.

% ============================================================
\section{Exercises}
\label{sec:ch07-exercises}
\begin{enumerate}[label=E7.\arabic*, leftmargin=*, itemsep=0.8em]
\item \textbf{Compute.} Find $\det A$ for $A=\begin{bmatrix}3&1&2\\0&-1&4\\1&0&1\end{bmatrix}$ by (a) row operations to triangular and (b) cofactor expansion along row 2. Compare.
\item \textbf{Invertibility.} For which $c$ is $A=\begin{bmatrix}1&c&0\\c&1&0\\0&0&2\end{bmatrix}$ singular? Relate to volume collapse and to eigenvalues.
\item \textbf{Rules.} Prove $\det(cA)=c^n\det A$ for $n\times n$ using the multilinear axiom. Use $\det(AB)=\det A\det B$ to show $\det(A^{-1})=1/\det A$ when invertible.
\item \textbf{Cramer.} Solve $\begin{cases}x+2y=5\\3x-y=1\end{cases}$ via Cramer's rule and compare with elimination. Why is Cramer $O(n!)$ and not used for $n>3$?
\item \textbf{Coding.} For random $n\times n$ matrices ($n=10,50,100$), compare \texttt{np.linalg.det} (via $LU$) with a cofactor implementation for $n\le6$ (time and overflow). Plot $\log|\det|$ vs.\ $n$.
\end{enumerate}
% End of Chapter 7

% SC1004 Chapter 8: Eigenvalues, Eigenvectors & Diagonalization (0->100 ground-up)
\chapter{Eigenvalues, Eigenvectors, and Diagonalization}
\label{ch:eigen}

\begin{quote}
\emph{``Eigenvectors are the directions a matrix leaves alone --- it only stretches them.''} Finding those directions turns a tangled linear map into independent scalings, powers into powers of scalars, and differential equations into decoupled ones.
\end{quote}

\noindent\fbox{\parbox{\dimexpr\linewidth-2\fboxsep-2\fboxrule}{%
\textbf{From zero: fixed directions before polynomials.} We start with ``which arrows stay on their own span?'' then derive $A\mathbf{v}=\lambda\mathbf{v}$, the characteristic polynomial, and diagonalization. No prior eigen theory assumed.}}
\vspace{0.4em}

\section*{Learning Outcomes}
After this chapter you will be able to:
\begin{enumerate}[label=(L\arabic*)]
    \item define eigenvalue/eigenvector and find them via $\det(A-\lambda I)=0$;
    \item compute eigenspaces and their dimensions (geometric multiplicity);
    \item test diagonalizability and construct $A=PDP^{-1}$ when possible;
    \item use diagonalization to compute $A^k$, recurrences, and matrix exponentials;
    \item recognise symmetric matrices' orthogonal diagonalization and the spectral picture.
\end{enumerate}

% ============================================================
\section{Eigenpairs: Directions That Stay Put}
\label{sec:eigen-def}
A nonzero $\mathbf{v}$ is an \emph{eigenvector} of $A\in\mathbb{R}^{n\times n}$ with \emph{eigenvalue} $\lambda\in\mathbb{R}$ (or $\mathbb{C}$) if
\[
  A\mathbf{v}=\lambda\mathbf{v}.
\]
Geometrically, $A$ stretches $\mathbf{v}$ by $\lambda$ but does not rotate it off its line. $\mathbf{0}$ is never an eigenvector (but $\lambda=0$ is allowed --- then $A\mathbf{v}=\mathbf{0}$).

\begin{figure}[h]
\centering
\begin{tikzpicture}[scale=0.85,every node/.style={font=\scriptsize}]
  \draw[->,thick] (-0.4,0)--(3.5,0) node[right] {$x$};
  \draw[->,thick] (0,-0.4)--(0,3.0) node[above] {$y$};
  \draw[->,very thick,blue!65!black] (0,0)--(2.0,0.6) node[midway,below] {$\mathbf{v}$};
  \draw[->,very thick,orange!75!black] (0,0)--(3.0,0.9) node[midway,above] {$A\mathbf{v}=1.5\mathbf{v}$};
  \draw[dashed,gray] (2.0,0.6)--(3.0,0.9);
  \node[font=\tiny,fill=white] at (1.5,1.4) {same line};
  \draw[->,very thick,red!65!black] (0,0)--(0.6,1.8) node[midway,left] {$\mathbf{w}$};
  \draw[->,very thick,violet!70!black] (0,0)--(1.8,1.2) node[midway,above] {$A\mathbf{w}$ not collinear};
  \draw[thick,red!60!black] (0.25,0.75)--(0.32,0.85)--(0.22,0.92);
  \node[font=\footnotesize\bfseries] at (1.5,-0.6) {$\mathbf{v}$ is eigen ($\lambda=1.5$); $\mathbf{w}$ is not};
\end{tikzpicture}
\caption{An eigenvector stays on its span under $A$; a generic vector is rotated off. The eigenvalue is the stretch factor (negative flips).}
\label{fig:eigen-direction}
\end{figure}

If $A\mathbf{v}=\lambda\mathbf{v}$ then $(A-\lambda I)\mathbf{v}=\mathbf{0}$ with $\mathbf{v}\neq\mathbf{0}$, so $A-\lambda I$ is singular. Hence

\begin{theorem}[Characteristic equation]
$\lambda$ is an eigenvalue of $A$ iff
\[
  p_A(\lambda)=\det(A-\lambda I_n)=0.
\]
$p_A$ is the \emph{characteristic polynomial} (degree $n$, leading term $(-\lambda)^n$ or $\lambda^n$ up to sign). Its roots are the eigenvalues; $\mathcal{N}(A-\lambda I)$ is the \emph{eigenspace} $E_\lambda$.
\end{theorem}

For $2\times2$, $A=\begin{bmatrix}a&b\\c&d\end{bmatrix}$,
\[
  p_A(\lambda)=\lambda^2-\operatorname{tr}(A)\lambda+\det(A),\quad \operatorname{tr}=a+d.
\]
So eigenvalues satisfy $\lambda_1+\lambda_2=\operatorname{tr}A$, $\lambda_1\lambda_2=\det A$ --- quick checks.

\begin{example}[A $2\times2$ shear vs.\ scaling]
$A=\begin{bmatrix}2&1\\0&3\end{bmatrix}$: $p_A(\lambda)=(2-\lambda)(3-\lambda)$, so $\lambda=2,3$. For $\lambda=2$, $(A-2I)=\begin{bmatrix}0&1\\0&1\end{bmatrix}$, $E_2=\operatorname{span}\{[1,0]^T\}$. For $\lambda=3$, $E_3=\operatorname{span}\{[1,1]^T\}$. Two independent eigen-directions --- diagonalizable (see \S\ref{sec:diag}).
\end{example}

\section{Eigenspaces and Multiplicities}
\label{sec:eigenspaces}
\emph{Algebraic multiplicity} $a_\lambda$: multiplicity of $\lambda$ as a root of $p_A$. \emph{Geometric multiplicity} $g_\lambda=\dim\mathcal{N}(A-\lambda I)$: number of independent eigenvectors for $\lambda$. Always $1\le g_\lambda\le a_\lambda$. Eigenvectors for distinct $\lambda$ are independent --- so $n$ distinct eigenvalues guarantees $n$ independent eigenvectors.

\begin{figure}[h]
\centering
\begin{tikzpicture}[scale=0.80,every node/.style={font=\scriptsize}]
  \draw[->,thick] (-0.5,0)--(4.5,0) node[right] {$\Re$};
  \draw[->,thick] (0,-1.5)--(0,2.0) node[above] {$\Im$};
  \fill[blue!60!black] (1.0,0) circle (2.5pt) node[below] {$\lambda_1=2$};
  \fill[red!60!black] (3.0,0) circle (2.5pt) node[below] {$\lambda_2=3$};
  \draw[dashed,gray] (1.0,0)--(1.0,1.4) node[above,font=\tiny,fill=white] {$E_{\lambda_1}$ line};
  \draw[dashed,gray] (3.0,0)--(3.0,1.4) node[above,font=\tiny,fill=white] {$E_{\lambda_2}$ line};
  \node[font=\tiny,align=center] at (2.0,-0.9) {Distinct $\lambda$ $\Rightarrow$ independent $\mathbf{v}$'s;\\eigenspaces intersect only at $\mathbf{0}$};
  \begin{scope}[xshift=6.5cm]
    \draw[->,thick] (-0.5,0)--(3.0,0) node[right] {$\lambda$};
    \node[font=\tiny,align=center] at (1.5,1.6) {Defective case:\\$a_\lambda=2$ but $g_\lambda=1$\\(only one eigen-line)};
    \draw[very thick,blue!60!black] (0,0)--(1.5,1.0) node[midway,above left] {$E_\lambda$};
    \fill[blue!60!black] (1.2,0) circle (2.5pt) node[below] {$\lambda$ (double root)};
    \node[font=\tiny,fill=orange!12,draw,rounded corners=2pt] at (1.5,-0.7) {not diagonalizable};
  \end{scope}
\end{tikzpicture}
\caption{Left: distinct eigenvalues give independent eigenspaces. Right: a repeated eigenvalue may have fewer eigenvectors than its algebraic multiplicity --- defective, not diagonalizable.}
\label{fig:eigenspaces-mult}
\end{figure}

\begin{example}[Defective]
$J=\begin{bmatrix}2&1\\0&2\end{bmatrix}$ has $p_J(\lambda)=(2-\lambda)^2$, so $a_2=2$, but $J-2I=\begin{bmatrix}0&1\\0&0\end{bmatrix}$ has nullity $1$, so $g_2=1<2$. Only one eigen-line --- $J$ is not diagonalizable (it is a Jordan block).
\end{example}

\section{Diagonalization}
\label{sec:diag}
$A$ is \emph{diagonalizable} if $A=PDP^{-1}$ with $D=\operatorname{diag}(\lambda_1,\dots,\lambda_n)$ and $P$ invertible. Columns of $P$ are then eigenvectors and diagonal entries are the corresponding eigenvalues.

\begin{theorem}[Diagonalization criterion]
$A\in\mathbb{R}^{n\times n}$ diagonalizable $\iff$ there are $n$ linearly independent eigenvectors $\iff$ for every $\lambda$, $g_\lambda=a_\lambda$ (equivalently $\sum_\lambda g_\lambda=n$).
\end{theorem}

Why care? Powers decouple: $A^k=PD^kP^{-1}=P\operatorname{diag}(\lambda_1^k,\dots,\lambda_n^k)P^{-1}$. So does $e^{At}=Pe^{Dt}P^{-1}$ and linear recurrences $\mathbf{x}_{k+1}=A\mathbf{x}_k$.

\begin{figure}[h]
\centering
\begin{tikzpicture}[scale=0.78,every node/.style={font=\scriptsize}]
  \node[draw,thick,rounded corners=3pt,fill=blue!10,minimum width=1.6cm,minimum height=0.9cm] (A) at (0,0) {$A$};
  \node[draw,thick,rounded corners=3pt,fill=orange!12,minimum width=1.2cm,minimum height=0.9cm] (Pinv) at (2.8,0.7) {$P^{-1}$};
  \node[draw,thick,rounded corners=3pt,fill=green!15,minimum width=1.6cm,minimum height=0.9cm] (D) at (5.0,0) {$D=\operatorname{diag}(\lambda_i)$};
  \node[draw,thick,rounded corners=3pt,fill=violet!12,minimum width=1.2cm,minimum height=0.9cm] (P) at (7.2,0.7) {$P$};
  \draw[->,very thick,blue!60!black] (A)--(Pinv) node[midway,above,sloped,font=\tiny] {change basis};
  \draw[->,very thick,green!55!black] (Pinv)--(D) node[midway,above,font=\tiny] {scale axes};
  \draw[->,very thick,violet!60!black] (D)--(P) node[midway,above,sloped,font=\tiny] {back};
  \node[font=\tiny,align=center] at (3.6,-0.9) {$A=PDP^{-1}$: in eigen-basis,\\$A$ acts as independent scalings $\lambda_i$};
  \draw[->,thick,gray,dashed] (A) to[out=-30,in=-150] (D) node[midway,below,font=\tiny] {$A^k=PD^kP^{-1}$};
\end{tikzpicture}
\caption{Diagonalization as change of basis: $P^{-1}$ moves to eigen-coordinates, $D$ scales each axis by $\lambda_i$, $P$ moves back. Powers act entrywise on $D$.}
\label{fig:diagonalization}
\end{figure}

\begin{example}[Fibonacci via diagonalization]
$\begin{bmatrix}F_{k+1}\\F_k\end{bmatrix}=\begin{bmatrix}1&1\\1&0\end{bmatrix}\begin{bmatrix}F_k\\F_{k-1}\end{bmatrix}$, so $\mathbf{x}_k=A^k\mathbf{x}_0$. $A$ has eigenvalues $\varphi=(1+\sqrt5)/2$, $\psi=(1-\sqrt5)/2$, $A=P\operatorname{diag}(\varphi,\psi)P^{-1}$, giving $F_k=(\varphi^k-\psi^k)/\sqrt5$.
\end{example}

\begin{example}[Python --- eigen and diagonalization]
\begin{lstlisting}[language=Python]
import numpy as np
A = np.array([[2.,1.],[0.,3.]])
w, V = np.linalg.eig(A)          # w eigenvalues, V columns are eigenvectors
print(w)                          # [2., 3.]
print(A @ V[:,0], w[0]*V[:,0])   # check A v = lambda v
# Diagonalization check:
P = V
D = np.diag(w)
print(np.allclose(A, P @ D @ np.linalg.inv(P)))
# Powers:
print(np.allclose(np.linalg.matrix_power(A, 5), P @ np.diag(w**5) @ np.linalg.inv(P)))
\end{lstlisting}
For symmetric $A$, \texttt{eigh} returns orthonormal $Q$ with $A=Q\Lambda Q^T$ (spectral theorem).
\end{example}

\section{Symmetric Matrices and the Spectral Theorem}
\label{sec:spectral}
If $A^T=A$ (real symmetric), then all eigenvalues are real, eigenvectors for distinct $\lambda$ are orthogonal, and $A$ is orthogonally diagonalizable: $A=Q\Lambda Q^T$ with $Q^TQ=I$, $\Lambda$ real diagonal. This is the \emph{spectral theorem} --- the basis that diagonalizes a symmetric matrix is a rotation/reflection. Applications: PCA (covariance is symmetric), graph Laplacians, quadratic forms $Q(\mathbf{x})=\mathbf{x}^TA\mathbf{x}=\sum\lambda_i y_i^2$.

% ============================================================

% --- Additional depth: Complex eigenvalues, Markov, and SVD preview ---
\section{Complex Eigenvalues and Real Blocks}
\label{sec:complex-eigen}
Real matrices can have complex eigenpairs when they rotate. $R_\theta=\begin{bmatrix}\cos\theta&-\sin\theta\\\sin\theta&\cos\theta\end{bmatrix}$ has $\lambda=e^{\pm i\theta}$, eigenvectors $[1,\mp i]^T$ in $\mathbb{C}^2$. Over $\mathbb{R}$, $R_\theta$ is not diagonalizable but block-diagonalizable with $2\times2$ rotation blocks. The real Schur form captures this; the complex Schur form is always triangular.

\section{Markov Chains and Steady States}
\label{sec:markov}
A column-stochastic $P$ (nonnegative, columns sum to $1$) has $\lambda=1$ as its largest eigenvalue (Perron--Frobenius); the eigenvector $\pi$ with $P\pi=\pi$ is the \emph{steady state}. Powers $P^k\mathbf{x}_0\to\pi$ when $P$ is primitive ($|\lambda_2|<1$). Example: PageRank is the $\lambda=1$ eigenvector of the web's link matrix.

\begin{example}[Two-state chain]
$P=\begin{bmatrix}0.9&0.2\\0.1&0.8\end{bmatrix}$: eigenvalues $1,0.7$, $E_1=\operatorname{span}\{[2,1]^T\}$, so $\pi=[2/3,1/3]^T$ after normalising to sum $1$. Starting from any distribution, $P^k\mathbf{x}_0\to\pi$ at rate $0.7^k$.
\end{example}

\section{Singular Value Decomposition}
\label{sec:svd}
Every $A\in\mathbb{R}^{m\times n}$ factors as $A=U\Sigma V^T$ with $U,V$ orthogonal and $\Sigma$ diagonal holding \emph{singular values} $\sigma_1\ge\sigma_2\ge\cdots\ge0$, where $\sigma_i=\sqrt{\lambda_i(A^TA)}$. Intuition: every linear map is rotate--stretch--rotate---$V^T$ aligns the input axes, $\Sigma$ stretches, $U$ orients the outputs. Unlike eigen-decomposition, SVD exists for every matrix, even rectangular ones.
\begin{example}[Worked $2\times2$]
$A=\begin{bmatrix}3&0\\0&-1\end{bmatrix}$ is diagonal in disguise: $A^TA=\operatorname{diag}(9,1)$, so $\sigma_1=3,\sigma_2=1$ with $V=I$ and $U=\operatorname{diag}(1,-1)$ absorbing the sign. The unit circle maps to an ellipse with semi-axes $3,1$---singular values are exactly the stretch factors.
\end{example}
The count of nonzero $\sigma_i$ equals $\operatorname{rank}A$. Dropping small $\sigma_i$ gives the best low-rank approximation (Eckart--Young): this is PCA---keep the top principal stretches, discard the noise. If $A=Q\Lambda Q^T$ is symmetric, $\sigma_i=|\lambda_i|$.



\begin{remark}[Diagonalizability vs.\ invertibility]
These are unrelated: $\begin{bmatrix}0&1\\0&0\end{bmatrix}$ is singular and not diagonalizable; $\begin{bmatrix}1&1\\0&1\end{bmatrix}$ is invertible but not diagonalizable; $I$ is both. Test diagonalizability via $g_\lambda=a_\lambda$, not via $\det$.
\end{remark}

\begin{example}[Quadratic forms]
For symmetric $A=Q\Lambda Q^T$, $f(\mathbf{x})=\mathbf{x}^TA\mathbf{x}=\sum\lambda_i y_i^2$ where $\mathbf{y}=Q^T\mathbf{x}$. So $\lambda_i>0$ for all $i$ $\iff$ $f(\mathbf{x})>0$ for $\mathbf{x}\neq\mathbf{0}$ (positive definite) --- Ch.~6's $A^TA$ is always positive semidefinite ($\lambda_i\ge0$).
\end{example}



\begin{example}[Power iteration intuition]
Repeatedly applying $A$ to any $\mathbf{x}_0$ with a component in the dominant eigenspace amplifies the largest $|\lambda|$ direction fastest: $A^k\mathbf{x}_0\approx c_1\lambda_1^k\mathbf{v}_1$. Normalising each step gives the \emph{power method} for $\lambda_1,\mathbf{v}_1$ --- how PageRank and PCA's top component are found when $n$ is huge and full \texttt{eig} is too costly.
\end{example}


\section{Chapter Notes}
Eigen (German ``own'') --- Hilbert's term for characteristic values. Diagonalization fails exactly on nontrivial Jordan blocks; over $\mathbb{C}$ every matrix is triangularizable (Schur) and Jordan-form. Numerically, eigenvalues are found via $QR$ iteration (Francis, 1961), not $\det(A-\lambda I)$.

% ============================================================
\section{Exercises}
\label{sec:ch08-exercises}
\begin{enumerate}[label=E8.\arabic*, leftmargin=*, itemsep=0.8em]
\item \textbf{Find eigenpairs.} Compute eigenvalues and eigenspaces for $A=\begin{bmatrix}4&2\\1&3\end{bmatrix}$ via $p_A(\lambda)=\lambda^2-7\lambda+10$. Give eigenvectors and verify $A\mathbf{v}=\lambda\mathbf{v}$.
\item \textbf{Diagonalize or not?} For $B=\begin{bmatrix}2&1&0\\0&2&0\\0&0&3\end{bmatrix}$, find $a_\lambda,g_\lambda$ for each $\lambda$. Is $B$ diagonalizable? If so give $P,D$; if not explain why.
\item \textbf{Powers.} Let $A=\begin{bmatrix}1&2\\2&1\end{bmatrix}=PDP^{-1}$ with $D=\operatorname{diag}(3,-1)$. Compute $A^4$ via $PD^4P^{-1}$ and verify by direct multiplication.
\item \textbf{Symmetric.} Show $S=\begin{bmatrix}2&1&1\\1&2&1\\1&1&2\end{bmatrix}$ is symmetric, find its eigenvalues ($4,1,1$) and an orthonormal eigenbasis. Write $S=Q\Lambda Q^T$ and $\mathbf{x}^TS\mathbf{x}$ in $y=Q^T\mathbf{x}$ coordinates.
\item \textbf{Coding.} For random symmetric $5\times5$ $A=(M+M^T)/2$, use \texttt{np.linalg.eigh} to get $Q,\Lambda$ and verify $Q^TQ\approx I$, $A\approx Q\Lambda Q^T$, and $\det A\approx\prod\lambda_i$, $\operatorname{tr}A\approx\sum\lambda_i$ numerically.
\end{enumerate}
% End of Chapter 8

\backmatter
\end{document}
